% !TEX root = ../main.tex

\chapter{مفاهیم نظری}

\section{مقدمه}

این فصل به معرفی مفاهیم نظری مورد نیاز برای درک طراحی و پیاده‌سازی
یک چت‌بات صوتی فارسی می‌پردازد.
سامانه مورد بررسی در این پروژه از چند مؤلفه اصلی شامل
\textbf{تشخیص خودکار گفتار}
\LTRfootnote{Automatic \lr{Speech} Recognition (\lr{ASR})}،
\textbf{مدل زبانی بزرگ}
\LTRfootnote{Large Language \lr{Model} (\lr{LLM})}
و
\textbf{تبدیل متن به گفتار}
\LTRfootnote{\lr{Text}-to-\lr{Speech} (\lr{TTS})}
تشکیل شده است.
علاوه بر این مؤلفه‌ها، پردازش‌هایی مانند تشخیص فعالیت صوتی،
مدیریت نویز، پس‌پردازش متن و حفظ بافت مکالمه نیز در عملکرد
یک سامانه صوتی تعاملی نقش مهمی دارند.

در ابتدا، مبانی پردازش گفتار و ساختار کلی سامانه‌های تشخیص گفتار
معرفی می‌شوند.
سپس معماری‌ها و رویکردهای اصلی مورد استفاده در مدل‌های
تشخیص گفتار بررسی خواهند شد؛ از جمله روش‌های مبتنی بر
\lr{wav2vec 2.0} و \lr{XLSR}، معماری \lr{Whisper}،
سامانه‌های مبتنی بر \lr{Kaldi/Vosk} و معماری‌های
\lr{Conformer} و \lr{FastConformer}.همچنین معیارهای متداول ارزیابی مدل‌های تشخیص گفتار،
از جمله نرخ خطای واژه
\LTRfootnote{\lr{Word Error Rate (WER)}}،
نرخ خطای نویسه
\LTRfootnote{\lr{Character Error Rate (CER)}}،
ضریب زمان واقعی
\LTRfootnote{\lr{Real-Time Factor (RTF)}}
و زمان تأخیر
\LTRfootnote{\lr{Latency}}
معرفی می‌شوند.

در ادامه، مفاهیم پایه مدل‌های زبانی بزرگ بررسی می‌شوند.
در این بخش، معماری \lr{Transformer}، مدل‌های مولد مبتنی بر
رمزگشا، مفهوم توکن‌سازی، فرآیند تولید متن و اهمیت مدیریت
بافت مکالمه در سامانه‌های گفت‌وگومحور تشریح خواهند شد.
با توجه به هدف پروژه مبنی بر اجرای محلی، موضوعاتی مانند
اندازه مدل، مصرف حافظه و سرعت استنتاج نیز مورد توجه قرار می‌گیرند.

پس از آن، مبانی سامانه‌های تبدیل متن به گفتار معرفی می‌شوند.
در این بخش، ساختار کلی سامانه‌های \lr{TTS}، معماری‌های مبتنی بر
\lr{VITS}، سامانه \lr{Piper} و معیارهای ارزیابی کیفیت گفتار
بررسی می‌شوند.
معیارهایی مانند
\textbf{میانگین امتیاز نظر}
\LTRfootnote{Mean Opinion Score (\lr{MOS})}
و شاخص‌های مرتبط با کیفیت ادراکی و قابل‌فهم‌بودن گفتار
در این قسمت معرفی خواهند شد.

در نهایت، مفاهیم مربوط به تشخیص فعالیت صوتی،
پردازش نویز و معماری یکپارچه سامانه‌های گفت‌وگوی صوتی
بررسی می‌شوند.
هدف این فصل ایجاد بنیان نظری لازم برای درک تصمیم‌های طراحی،
انتخاب مدل‌ها و روش ارزیابی سامانه‌ای است که در فصل‌های بعدی
تشریح خواهد شد.
% =========================================================
% 2-2
% =========================================================
\section{مبانی پردازش گفتار}

گفتار انسان یک سیگنال صوتی پیوسته در زمان است که در اثر ارتعاش
تارهای صوتی و تغییر شکل مجرای گفتاری تولید می‌شود.
برای پردازش این سیگنال توسط سامانه‌های رایانه‌ای، ابتدا باید صوت
آنالوگ به یک نمایش دیجیتال تبدیل شود.
نمایش دیجیتال گفتار پایه بسیاری از وظایف پردازش صوت، از جمله
تشخیص گفتار، تشخیص فعالیت صوتی، کاهش نویز و تبدیل گفتار به
بازنمایی‌های قابل استفاده برای مدل‌های یادگیری ماشین است.

در یک سامانه گفت‌وگوی صوتی، سیگنال خام دریافت‌شده از میکروفون
معمولاً مستقیماً در اختیار مدل قرار نمی‌گیرد.
پیش از آن، عملیات‌هایی نظیر نمونه‌برداری، تبدیل قالب داده،
تقسیم سیگنال به بازه‌های کوتاه زمانی و در برخی سامانه‌ها
استخراج نمایش‌های حوزه فرکانس یا ویژگی‌های صوتی انجام می‌شود.
اگرچه بسیاری از مدل‌های جدید قادرند بخش زیادی از فرآیند استخراج
ویژگی را به‌صورت خودکار بیاموزند، آشنایی با مفاهیم پایه پردازش
گفتار برای درک نحوه عملکرد این مدل‌ها و تفسیر نتایج آن‌ها ضروری است.


% ---------------------------------------------------------
\subsection{نمونه‌برداری و نمایش دیجیتال گفتار}

سیگنال صوتی در محیط فیزیکی یک سیگنال پیوسته است.
برای ذخیره و پردازش آن در رایانه، دامنه سیگنال در بازه‌های زمانی
مشخص اندازه‌گیری می‌شود که این فرآیند
\textbf{نمونه‌برداری}
\LTRfootnote{Sampling}
نام دارد.

اگر سیگنال پیوسته صوت را با $x(t)$ نمایش دهیم، نمایش گسسته آن
را می‌توان به‌صورت زیر نوشت:

\begin{equation}
x[n] = x(nT_s)
\end{equation}

که در آن $T_s$ فاصله زمانی میان دو نمونه متوالی است.
نرخ نمونه‌برداری
\LTRfootnote{Sampling Rate}
برابر است با:

\begin{equation}
f_s = \frac{1}{T_s}
\end{equation}

و نشان می‌دهد که در هر ثانیه چند نمونه از سیگنال صوتی ثبت
می‌شود.

بر اساس قضیه نمونه‌برداری نایکوئیست--شنون، برای بازسازی مناسب
یک سیگنال، نرخ نمونه‌برداری باید حداقل دو برابر بیشترین فرکانس
موجود در سیگنال باشد.
در کاربردهای مرتبط با تشخیص گفتار، نرخ‌های نمونه‌برداری
\lr{16 kHz} بسیار متداول هستند، زیرا بخش عمده اطلاعات مورد نیاز
برای تشخیص محتوای گفتار در این بازه فرکانسی قرار دارد.

علاوه بر نرخ نمونه‌برداری، تعداد بیت‌های مورد استفاده برای نمایش
دامنه هر نمونه نیز بر دقت نمایش دیجیتال اثر می‌گذارد.
برای مثال، نمایش‌های \lr{16-bit PCM}
یکی از قالب‌های رایج برای ذخیره صوت گفتاری هستند.

در یک خط لوله تشخیص گفتار، هماهنگ‌بودن نرخ نمونه‌برداری ورودی
با نرخ مورد انتظار مدل اهمیت زیادی دارد.
در صورتی که صوت با نرخ متفاوتی ضبط شده باشد، معمولاً لازم است
عملیات
\textbf{بازنمونه‌برداری}
\LTRfootnote{Resampling}
پیش از اجرای مدل انجام شود.


% ---------------------------------------------------------
\subsection{قاب‌بندی و پنجره‌بندی سیگنال}

سیگنال گفتار در بازه‌های زمانی طولانی یک سیگنال غیرایستا
\LTRfootnote{Non-stationary}
محسوب می‌شود؛ زیرا ویژگی‌های آن در طول زمان تغییر می‌کنند.
با این حال، در بازه‌های زمانی بسیار کوتاه می‌توان رفتار گفتار
را تقریباً ایستا در نظر گرفت.

به همین دلیل، یکی از روش‌های متداول در پردازش گفتار تقسیم
سیگنال به قطعات کوتاه موسوم به
\textbf{قاب}
\LTRfootnote{Frame}
است.
طول هر قاب معمولاً در محدوده چند ده میلی‌ثانیه انتخاب می‌شود.
قاب‌های متوالی اغلب با یکدیگر هم‌پوشانی دارند تا از حذف تغییرات
مهم در مرز میان قاب‌ها جلوگیری شود.

برای کاهش ناپیوستگی در ابتدا و انتهای هر قاب، از یک تابع
\textbf{پنجره}
\LTRfootnote{Window Function}
استفاده می‌شود.
یکی از نمونه‌های رایج، پنجره همینگ
\LTRfootnote{Hamming Window}
است.

اگر $x[n]$ یک قاب از سیگنال و $w[n]$ تابع پنجره باشد،
سیگنال پنجره‌شده به شکل زیر تعریف می‌شود:

\begin{equation}
x_w[n] = x[n]w[n]
\end{equation}

قاب‌بندی و پنجره‌بندی به‌ویژه در روش‌هایی که از نمایش
زمان--فرکانس استفاده می‌کنند اهمیت دارند.


% ---------------------------------------------------------
\subsection{نمایش حوزه زمان و حوزه فرکانس}

ساده‌ترین نمایش یک سیگنال گفتار، نمایش تغییرات دامنه آن در
طول زمان است.
این نمایش که
\textbf{شکل موج}
\LTRfootnote{\lr{Waveform}}
نامیده می‌شود، اطلاعاتی درباره دامنه، سکوت، شدت تقریبی و
ساختار زمانی سیگنال ارائه می‌دهد.

با این حال، بسیاری از ویژگی‌های مهم گفتار در حوزه فرکانس
قابل مشاهده‌تر هستند.
برای تبدیل سیگنال از حوزه زمان به حوزه فرکانس می‌توان از
تبدیل فوریه استفاده کرد.
در پردازش گفتار، به دلیل تغییر ویژگی‌های سیگنال در طول زمان،
معمولاً از
\textbf{تبدیل فوریه کوتاه‌مدت}
\LTRfootnote{Short-Time Fourier Transform (STFT)}
استفاده می‌شود.

برای یک سیگنال $x[n]$، تبدیل فوریه کوتاه‌مدت را می‌توان به
شکل کلی زیر بیان کرد:

\begin{equation}
X(m,\omega)
=
\sum_{n=-\infty}^{+\infty}
x[n]\,w[n-m]\,e^{-j\omega n}
\end{equation}

که در آن $w[n]$ تابع پنجره و $m$ موقعیت زمانی قاب است.

با محاسبه اندازه ضرایب تبدیل فوریه برای قاب‌های متوالی،
می‌توان نمایش دوبعدی زمان--فرکانس موسوم به
\textbf{طیف‌نگاشت}
\LTRfootnote{\lr{Spectrogram}}
را ایجاد کرد.
در طیف‌نگاشت، یک محور زمان، محور دیگر فرکانس و مقدار هر نقطه
معمولاً نشان‌دهنده شدت یا انرژی مؤلفه فرکانسی مربوطه است.

طیف‌نگاشت‌ها نقش مهمی در پردازش گفتار دارند و مبنای بسیاری از
نمایش‌های ورودی مورد استفاده در سامانه‌های تشخیص گفتار و
مدل‌های صوتی هستند.


% ---------------------------------------------------------
\subsection{مقیاس Mel و ویژگی‌های گفتاری}

ادراک انسان از فرکانس به‌صورت کاملاً خطی نیست.
تغییرات فرکانسی در نواحی فرکانس پایین با حساسیت بیشتری نسبت
به تغییرات مشابه در فرکانس‌های بالا درک می‌شوند.
برای نزدیک‌کردن نمایش فرکانسی به نحوه ادراک انسان، از
\textbf{مقیاس مل}
\LTRfootnote{Mel Scale}
استفاده می‌شود.

یک رابطه رایج برای تبدیل فرکانس بر حسب هرتز به مقیاس مل
به‌صورت زیر است:

\begin{equation}
m
=
2595
\log_{10}
\left(
1+\frac{f}{700}
\right)
\end{equation}

که در آن $f$ فرکانس بر حسب هرتز و $m$ مقدار متناظر در مقیاس
مل است.

یکی از نمایش‌های رایج مبتنی بر این مفهوم،
\textbf{طیف‌نگاشت Mel}
\LTRfootnote{Mel \lr{Spectrogram}}
است.
در این نمایش، انرژی طیفی سیگنال پس از اعمال مجموعه‌ای از
فیلترهای قرارگرفته بر روی مقیاس مل محاسبه می‌شود.

در سامانه‌های سنتی‌تر تشخیص گفتار، ویژگی‌هایی مانند
\textbf{ضرایب کپسترال فرکانس Mel}
\LTRfootnote{Mel-Frequency Cepstral Coefficients (MFCC)}
به‌طور گسترده مورد استفاده قرار گرفته‌اند.
این ویژگی‌ها تلاش می‌کنند اطلاعات مهم طیفی گفتار را در یک
بردار با ابعاد کمتر نمایش دهند.

در مقابل، بسیاری از مدل‌های یادگیری عمیق جدید وابستگی کمتری
به استخراج دستی ویژگی دارند.
برخی مدل‌ها مستقیماً شکل موج خام را دریافت کرده و بازنمایی
مناسب را طی فرآیند آموزش استخراج می‌کنند، در حالی که برخی
دیگر از طیف‌نگاشت یا ویژگی‌های مبتنی بر فیلترهای Mel به‌عنوان
ورودی استفاده می‌کنند.


% ---------------------------------------------------------
\subsection{نویز و نسبت سیگنال به نویز}

در شرایط واقعی، صوت دریافت‌شده از میکروفون معمولاً تنها شامل
گفتار کاربر نیست و ممکن است نویزهای مختلف محیطی نیز در آن
وجود داشته باشند.
صدای سیستم تهویه، خیابان، فن رایانه، موسیقی پس‌زمینه یا
گفتار افراد دیگر نمونه‌هایی از عواملی هستند که می‌توانند بر
کیفیت ورودی اثر بگذارند.

یکی از معیارهای متداول برای توصیف میزان نویز موجود در سیگنال،
\textbf{نسبت سیگنال به نویز}
\LTRfootnote{Signal-to-\lr{Noise} Ratio (SNR)}
است.
این معیار معمولاً بر حسب دسی‌بل بیان می‌شود:

\begin{equation}
\mathrm{SNR}_{dB}
=
10
\log_{10}
\left(
\frac{P_{\mathrm{signal}}}
     {P_{\mathrm{noise}}}
\right)
\end{equation}

که در آن $P_{\mathrm{signal}}$ توان سیگنال اصلی و
$P_{\mathrm{noise}}$ توان نویز است.

هرچه مقدار \lr{SNR} بیشتر باشد، سهم گفتار نسبت به نویز بالاتر
است.
کاهش \lr{SNR} می‌تواند موجب افزایش خطای سامانه تشخیص گفتار شود،
به‌ویژه زمانی که نویز دارای مؤلفه‌های فرکانسی مشابه گفتار باشد.

در یک سامانه مکالمه‌ای واقعی، بررسی عملکرد مدل تشخیص گفتار
فقط بر روی فایل‌های صوتی تمیز کافی نیست.
ازاین‌رو، مقاومت سامانه در برابر نویز یکی از موضوعاتی است که
می‌تواند در مرحله ارزیابی تجربی مورد بررسی قرار گیرد.


% ---------------------------------------------------------
\subsection{سکوت و فعالیت صوتی}

جریان صوتی دریافت‌شده از میکروفون علاوه بر گفتار شامل
بازه‌های سکوت نیز است.
ارسال تمام این جریان بدون تشخیص زمان واقعی گفتار می‌تواند
باعث افزایش پردازش غیرضروری و ایجاد مشکلاتی در تعیین زمان
پایان جمله کاربر شود.

به همین دلیل، در بسیاری از سامانه‌های صوتی از
\textbf{تشخیص فعالیت صوتی}
\LTRfootnote{Voice Activity Detection (VAD)}
استفاده می‌شود.
وظیفه یک سامانه \lr{VAD} تعیین این است که در هر بازه زمانی
آیا گفتار انسانی وجود دارد یا خیر.

به‌صورت ساده می‌توان تصمیم یک سامانه تشخیص فعالیت صوتی را
به شکل زیر نمایش داد:

\begin{equation}
V(t)
=
\begin{cases}
1 & \mbox{\rl{وجود گفتار}} \\
0 & \mbox{\rl{عدم وجود گفتار}}
\end{cases}
\end{equation}

تشخیص دقیق شروع و پایان گفتار در یک چت‌بات صوتی اهمیت زیادی
دارد، زیرا خروجی این بخش می‌تواند تعیین کند که چه زمانی نمونه
صوتی کامل شده و برای مدل تشخیص گفتار ارسال شود.
خطای این مؤلفه ممکن است باعث حذف ابتدای یا انتهای جمله و یا
ارسال بخش‌های طولانی سکوت به مدل شود.


% ---------------------------------------------------------
\subsection{نقش پیش‌پردازش در خط لوله گفتار}

هدف از پیش‌پردازش صوت ایجاد ورودی مناسب‌تر و پایدارتر برای
مراحل بعدی سامانه است.
بسته به نوع مدل و شرایط اجرا، این پیش‌پردازش می‌تواند شامل
تبدیل نرخ نمونه‌برداری، تبدیل صوت چندکاناله به تک‌کاناله،
نرمال‌سازی دامنه، حذف سکوت‌های اضافی، تشخیص فعالیت صوتی و
کاهش نویز باشد.

با این حال، اعمال پیش‌پردازش بیشتر لزوماً به معنای افزایش
دقت نیست.
مدل‌های مدرن تشخیص گفتار ممکن است در طول آموزش با تنوع زیادی
از شرایط صوتی مواجه شده باشند و برخی عملیات پردازشی شدید
می‌تواند بخشی از اطلاعات مفید سیگنال را نیز حذف کند.
ازاین‌رو، تأثیر هر مرحله پیش‌پردازش باید به‌صورت تجربی
بررسی شود.

در معماری مورد بررسی در این پروژه، پردازش صوت در ابتدای
خط لوله قرار دارد و خروجی آن به بخش تشخیص خودکار گفتار
ارسال می‌شود.
مفاهیم معرفی‌شده در این بخش پایه لازم برای بررسی روش‌ها و
معماری‌های تشخیص گفتار در بخش بعدی را فراهم می‌کنند.

% =========================================================
% 2-3
% =========================================================
\section{تشخیص خودکار گفتار}

\textbf{تشخیص خودکار گفتار}
\LTRfootnote{Automatic \lr{Speech} Recognition (\lr{ASR})}
فرآیندی است که در آن یک سامانه رایانه‌ای سیگنال گفتاری را دریافت
کرده و محتوای زبانی آن را به متن تبدیل می‌کند.
به‌صورت کلی، اگر دنباله نمونه‌های صوتی ورودی را با $X$ و متن
متناظر را با $Y$ نمایش دهیم، هدف یک سامانه تشخیص گفتار یافتن
محتمل‌ترین دنباله متنی برای صوت ورودی است:

\begin{equation}
\hat{Y}
=
\arg\max_Y P(Y|X)
\end{equation}

در این رابطه، $P(Y|X)$ احتمال تولید دنباله متنی $Y$ با داشتن
سیگنال صوتی $X$ است و $\hat{Y}$ متن نهایی پیش‌بینی‌شده توسط
سامانه را نشان می‌دهد.

تشخیص گفتار مسئله‌ای پیچیده است؛ زیرا یک جمله واحد می‌تواند
توسط افراد مختلف با سرعت، لهجه، شدت صدا و ویژگی‌های آکوستیکی
متفاوت بیان شود.
علاوه بر این، وجود نویز محیط، کیفیت میکروفون، تفاوت در تلفظ،
هم‌آوایی برخی واژه‌ها و وابستگی معنای کلمات به بافت جمله می‌تواند
فرآیند تبدیل گفتار به متن را دشوارتر کند.

سامانه‌های تشخیص گفتار در طول زمان از معماری‌های مبتنی بر
مدل‌های آماری و اجزای جداگانه به سمت سامانه‌های
\textbf{انتهابه‌انتها}
\LTRfootnote{End-to-End}
حرکت کرده‌اند.
در سامانه‌های جدید، شبکه‌های عصبی عمیق می‌توانند بخش قابل توجهی
از نگاشت میان صوت ورودی و دنباله متنی را مستقیماً یاد بگیرند.
مدل‌هایی مانند \lr{wav2vec 2.0}، \lr{Whisper} و
معماری‌های مبتنی بر \lr{Conformer} نمونه‌هایی از این روند هستند
\cite{baevski2020wav2vec2,radford2022whisper,gulati2020conformer}.


% ---------------------------------------------------------
\subsection{ساختار کلی یک سامانه تشخیص گفتار}

یک سامانه تشخیص گفتار را می‌توان به‌صورت مفهومی شامل چند مرحله
در نظر گرفت.
ابتدا سیگنال صوتی دریافت و به شکل مناسب برای مدل آماده می‌شود.
این مرحله می‌تواند شامل تغییر نرخ نمونه‌برداری، نرمال‌سازی صوت،
حذف بخش‌های غیرضروری و استخراج ویژگی‌های صوتی باشد.

پس از آماده‌سازی صوت، بخش اصلی مدل تلاش می‌کند دنباله‌ای از
بازنمایی‌های صوتی را به واحدهای زبانی مرتبط سازد.
این واحدها ممکن است نویسه، زیرواژه
\LTRfootnote{Subword}
یا توکن باشند.

در مرحله نهایی، فرآیند
\textbf{رمزگشایی}
\LTRfootnote{Decoding}
برای انتخاب دنباله متنی نهایی انجام می‌شود.
نوع رمزگشایی بسته به معماری مدل متفاوت است.
برخی مدل‌ها از پیش‌بینی مستقل یا تقریباً مستقل برچسب‌های زمانی
استفاده می‌کنند، در حالی که برخی دیگر متن را به‌صورت
خودبازگشتی
\LTRfootnote{Autoregressive}
توکن به توکن تولید می‌کنند.

به‌صورت کلی می‌توان جریان پردازش را به شکل زیر نمایش داد:

\begin{center}
صوت ورودی
$\longrightarrow$
پیش‌پردازش
$\longrightarrow$
استخراج بازنمایی
$\longrightarrow$
مدل تشخیص گفتار
$\longrightarrow$
رمزگشایی
$\longrightarrow$
متن
\end{center}

در سامانه پیشنهادی این پروژه، خروجی بخش تشخیص گفتار مستقیماً
به بخش پردازش متنی و سپس مدل زبانی ارسال می‌شود.
بنابراین کیفیت خروجی این مرحله تأثیر مستقیمی بر عملکرد مراحل
بعدی دارد.


% ---------------------------------------------------------
\subsection{روش CTC}

یکی از روش‌های مهم مورد استفاده در مدل‌های تشخیص گفتار،
\textbf{طبقه‌بندی زمانی پیوندگرا}
\LTRfootnote{Connectionist Temporal Classification (CTC)}
است.

مسئله اصلی در بسیاری از مجموعه‌داده‌های گفتاری این است که
متن صحیح جمله در دسترس است، اما مشخص نیست هر نویسه یا توکن
دقیقاً به کدام بخش زمانی از سیگنال صوت مربوط می‌شود.
روش \lr{CTC} امکان آموزش مدل را بدون نیاز به تعیین صریح
هم‌ترازی میان فریم‌های صوت و برچسب‌های متنی فراهم می‌کند.

در این روش، یک نماد ویژه موسوم به
\textbf{blank}
به فضای خروجی اضافه می‌شود.
مدل می‌تواند در هر گام زمانی یک نویسه یا نماد خالی را پیش‌بینی
کند.
سپس با حذف نمادهای خالی و ادغام پیش‌بینی‌های تکراری،
دنباله نهایی متن تشکیل می‌شود.

برای مثال، دنباله‌ای مانند:

\begin{equation}
\_ \; \mbox{\rl{س}} \; \mbox{\rl{س}} \; \_ \; \mbox{\rl{ل}} \; \_ \; \mbox{\rl{ا}} \; \mbox{\rl{ا}} \; \mbox{\rl{م}}
\end{equation}

می‌تواند پس از اعمال قوانین \lr{CTC} به واژه:

\begin{center}
سلام
\end{center}

تبدیل شود.

احتمال یک متن نهایی در \lr{CTC} از مجموع احتمال تمام مسیرهای
ممکنی محاسبه می‌شود که پس از حذف نمادهای خالی و تکرارها به
آن متن منجر می‌شوند.

این رویکرد در بسیاری از مدل‌های تشخیص گفتار انتهابه‌انتها
به‌کار گرفته شده است و مدل‌های مبتنی بر \lr{wav2vec 2.0}
نیز می‌توانند با استفاده از هدف \lr{CTC} برای تشخیص گفتار
تنظیم شوند \cite{baevski2020wav2vec2}.


% ---------------------------------------------------------
\subsection{مدل‌های دنباله‌به‌دنباله و رمزگشایی خودبازگشتی}

رویکرد دیگر در تشخیص گفتار استفاده از مدل‌های
\textbf{دنباله‌به‌دنباله}
\LTRfootnote{Sequence-to-Sequence}
است.
در این معماری‌ها، مدل ابتدا بازنمایی صوت ورودی را پردازش
می‌کند و سپس متن را به‌صورت یک دنباله از توکن‌ها تولید می‌کند.

در مدل‌های خودبازگشتی، احتمال دنباله خروجی را می‌توان به شکل
زیر تجزیه کرد:

\begin{equation}
P(Y|X)
=
\prod_{t=1}^{T}
P(y_t | y_1,\ldots,y_{t-1},X)
\end{equation}

در این رابطه، تولید هر توکن جدید نه‌تنها به صوت ورودی، بلکه
به توکن‌های تولیدشده قبلی نیز وابسته است.

این ویژگی می‌تواند به مدل کمک کند تا از اطلاعات زبانی موجود
در دنباله قبلی برای پیش‌بینی توکن بعدی استفاده کند.
در مقابل، تولید خودبازگشتی معمولاً به اجرای چندین مرحله متوالی
نیاز دارد و می‌تواند نسبت به برخی روش‌های غیرخودبازگشتی
هزینه زمانی بیشتری داشته باشد.

مدل \lr{Whisper} نمونه‌ای از یک سامانه مبتنی بر معماری
\lr{Transformer} و رمزگشایی دنباله‌به‌دنباله است که برای
وظایف مختلف مرتبط با گفتار آموزش داده شده است
\cite{radford2022whisper}.


% ---------------------------------------------------------
\subsection{رمزگشایی Greedy و Beam Search}

پس از محاسبه احتمال توکن‌های خروجی، سامانه باید دنباله نهایی
متن را انتخاب کند.
یکی از ساده‌ترین روش‌ها،
\textbf{رمزگشایی حریصانه}
\LTRfootnote{Greedy Decoding}
است.

در این روش، در هر گام توکنی انتخاب می‌شود که بیشترین احتمال
را در همان مرحله دارد:

\begin{equation}
\hat{y}_t
=
\arg\max_y P(y|X)
\end{equation}

این روش ساده و سریع است، اما تضمین نمی‌کند که دنباله نهایی
بهترین دنباله ممکن باشد.

روش دیگر
\textbf{جستجوی پرتو}
\LTRfootnote{Beam Search}
است.
در این روش به‌جای نگهداری تنها یک انتخاب در هر مرحله، چندین
دنباله با بیشترین احتمال حفظ می‌شوند.
تعداد این دنباله‌ها با پارامتر
\textbf{عرض پرتو}
\LTRfootnote{Beam Width}
تعیین می‌شود.

افزایش عرض پرتو می‌تواند در برخی شرایط کیفیت رمزگشایی را
بهبود دهد، اما هزینه محاسباتی و زمان استنتاج را نیز افزایش
می‌دهد.
بنابراین، نوع رمزگشایی می‌تواند بر توازن میان سرعت و دقت
سامانه تأثیر بگذارد.


% ---------------------------------------------------------
\subsection{نرخ خطای واژه}

یکی از مهم‌ترین معیارهای استاندارد برای ارزیابی سامانه‌های
تشخیص گفتار،
\textbf{نرخ خطای واژه}
\LTRfootnote{Word Error Rate (\lr{WER})}
است.

برای محاسبه این معیار، متن تولیدشده توسط مدل با متن مرجع
مقایسه می‌شود.
سه نوع خطا در نظر گرفته می‌شوند:

\begin{itemize}

    \item
    \textbf{جایگزینی}
    \LTRfootnote{Substitution}:
    یک واژه با واژه دیگری اشتباه تشخیص داده شده است.

    \item
    \textbf{حذف}
    \LTRfootnote{Deletion}:
    یک واژه موجود در متن مرجع در خروجی مدل وجود ندارد.

    \item
    \textbf{درج}
    \LTRfootnote{Insertion}:
    مدل واژه‌ای تولید کرده است که در متن مرجع وجود ندارد.

\end{itemize}

اگر تعداد جایگزینی‌ها را با $S$، حذف‌ها را با $D$،
درج‌ها را با $I$ و تعداد کل واژه‌های متن مرجع را با $N$
نمایش دهیم، نرخ خطای واژه به‌صورت زیر تعریف می‌شود:

\begin{equation}
WER
=
\frac{S+D+I}{N}
\end{equation}

مقدار کمتر \lr{WER} نشان‌دهنده عملکرد بهتر مدل است.
برای مثال، مقدار:

\begin{equation}
WER = 0.10
\end{equation}

به این معنا است که نسبت تعداد عملیات خطا به تعداد واژه‌های
مرجع برابر با $10\%$ بوده است.

با این حال، تفسیر \lr{WER} باید با توجه به زبان،
نوع داده، روش نرمال‌سازی متن و مجموعه آزمایش انجام شود.
برای زبان فارسی، مواردی مانند فاصله و نیم‌فاصله، شکل نوشتاری
برخی کلمات و نحوه نرمال‌سازی متن می‌توانند بر مقدار نهایی
این معیار اثرگذار باشند.


% ---------------------------------------------------------
\subsection{نرخ خطای نویسه}

معیار دیگر،
\textbf{نرخ خطای نویسه}
\LTRfootnote{Character Error Rate (\lr{CER})}
است.
این معیار مشابه \lr{WER} محاسبه می‌شود، با این تفاوت که
واحد مقایسه به‌جای واژه، نویسه است:

\begin{equation}
CER
=
\frac{S_c+D_c+I_c}{N_c}
\end{equation}

که در آن $S_c$، $D_c$ و $I_c$ به‌ترتیب تعداد جایگزینی،
حذف و درج نویسه‌ها و $N_c$ تعداد نویسه‌های متن مرجع است.

استفاده هم‌زمان از \lr{WER} و \lr{CER} در ارزیابی زبان فارسی
می‌تواند مفید باشد.
ممکن است یک مدل تنها بخشی از یک واژه را اشتباه تشخیص دهد،
که در سطح واژه یک خطای کامل محسوب می‌شود، اما در سطح نویسه
میزان اختلاف کوچک‌تر باشد.

ازاین‌رو، در ارزیابی مدل‌های تشخیص گفتار این پروژه هر دو معیار
\lr{WER} و \lr{CER} مورد استفاده قرار گرفته‌اند.


% ---------------------------------------------------------
\subsection{ضریب زمان واقعی}

در یک چت‌بات صوتی، دقت تنها معیار مهم نیست.
مدل باید بتواند صوت را با سرعت کافی پردازش کند تا تعامل با
کاربر با تأخیر زیاد همراه نباشد.

یکی از معیارهای متداول برای اندازه‌گیری سرعت پردازش صوت،
\textbf{ضریب زمان واقعی}
\LTRfootnote{Real-Time Factor (\lr{RTF})}
است که به شکل زیر تعریف می‌شود:

\begin{equation}
RTF
=
\frac{T_{\mathrm{inference}}}
     {T_{\mathrm{audio}}}
\end{equation}

که در آن $T_{\mathrm{inference}}$ زمان مورد نیاز برای پردازش
و $T_{\mathrm{audio}}$ مدت زمان فایل صوتی است.

برای مثال، اگر پردازش یک فایل صوتی ده‌ثانیه‌ای دو ثانیه
زمان ببرد:

\begin{equation}
RTF
=
\frac{2}{10}
=
0.2
\end{equation}

خواهد بود.

مقدار \lr{RTF}<1 به این معنا است که مدل صوت را سریع‌تر از
زمان واقعی پردازش می‌کند.
هرچه مقدار \lr{RTF} کمتر باشد، سرعت پردازش نسبت به طول صوت
بیشتر است.

این معیار برای پروژه حاضر اهمیت ویژه‌ای دارد؛ زیرا هدف تنها
دستیابی به کمترین نرخ خطا نیست، بلکه مدل منتخب باید برای
استفاده در یک سامانه مکالمه‌ای محلی از سرعت کافی نیز برخوردار
باشد.


% ---------------------------------------------------------
\subsection{زمان تأخیر}

معیار دیگری که در این پروژه مورد استفاده قرار می‌گیرد،
\textbf{زمان تأخیر}
\LTRfootnote{\lr{Latency}}
است.

در ارزیابی مستقل یک مدل تشخیص گفتار، زمان تأخیر را می‌توان
به‌صورت مدت زمان مورد نیاز برای پردازش یک نمونه صوتی تعریف کرد.

اگر تعداد نمونه‌های ارزیابی برابر با $M$ و زمان استنتاج نمونه
$i$ام برابر با $t_i$ باشد، متوسط زمان تأخیر به شکل زیر است:

\begin{equation}
\mathrm{Average\ Latency}
=
\frac{1}{M}
\sum_{i=1}^{M} t_i
\end{equation}

این معیار با \lr{RTF} متفاوت است.
\lr{RTF} زمان پردازش را نسبت به طول صوت نرمال می‌کند، در حالی
که متوسط تأخیر مدت زمان واقعی انتظار برای پردازش هر نمونه را
نشان می‌دهد.

در سامانه نهایی، مفهوم تأخیر گسترده‌تر است و علاوه بر
\lr{ASR}، زمان پردازش مدل زبانی، سامانه تبدیل متن به گفتار
و سایر پردازش‌های جانبی نیز در تأخیر انتهابه‌انتها نقش دارند.


% ---------------------------------------------------------
\subsection{مصرف منابع محاسباتی}

اجرای محلی مدل‌های تشخیص گفتار نیازمند توجه به میزان مصرف
منابع سخت‌افزاری است.
مدل‌هایی با تعداد پارامترهای زیاد ممکن است دقت مناسبی ارائه دهند،
اما برای اجرا به حافظه گرافیکی یا قدرت پردازشی بیشتری نیاز
داشته باشند.

در این پروژه، علاوه بر معیارهای دقت و سرعت،
میزان مصرف حافظه
\lr{GPU}
نیز در مقایسه مدل‌ها در نظر گرفته شده است.

این موضوع در کاربردهای محلی اهمیت زیادی دارد؛ زیرا مدل منتخب
باید در کنار مدل زبانی و سامانه تبدیل متن به گفتار بر روی
یک سیستم واحد اجرا شود.
بنابراین استفاده تمام حافظه گرافیکی توسط بخش \lr{ASR} می‌تواند
امکان اجرای هم‌زمان سایر اجزای سامانه را محدود کند.

در نتیجه، انتخاب مدل مناسب یک مسئله چندمعیاره است و باید
توازن میان موارد زیر برقرار شود:

\begin{itemize}
    \item دقت تشخیص گفتار؛
    \item سرعت پردازش؛
    \item زمان تأخیر؛
    \item مصرف حافظه و منابع محاسباتی؛
    \item قابلیت اجرای محلی؛
    \item و سازگاری با سایر مؤلفه‌های خط لوله.
\end{itemize}

معیارهای معرفی‌شده در این بخش، مبنای ارزیابی تجربی مدل‌های
تشخیص گفتار در این پروژه هستند.
در فصل آزمایش‌ها، مقادیر \lr{WER}، \lr{CER}، \lr{RTF}،
زمان تأخیر و مصرف حافظه برای مدل‌های مختلف گزارش و با یکدیگر
مقایسه خواهند شد.
% =========================================================
% 2-4
% =========================================================
\section{معماری‌های نوین تشخیص گفتار}

در سال‌های اخیر، معماری سامانه‌های تشخیص گفتار تغییرات
قابل توجهی داشته است.
روش‌های سنتی معمولاً از چند مؤلفه مستقل مانند مدل آکوستیکی،
مدل زبانی و الگوریتم رمزگشایی تشکیل می‌شدند، در حالی که
بسیاری از سامانه‌های جدید تلاش می‌کنند بخش بیشتری از فرآیند
نگاشت صوت به متن را به‌صورت انتهابه‌انتها یاد بگیرند.

در این پروژه، چند خانواده متفاوت از مدل‌های تشخیص گفتار
برای زبان فارسی مورد ارزیابی قرار گرفته‌اند.
این مدل‌ها شامل یک مدل مبتنی بر
\lr{wav2vec 2.0/XLSR}،
مدل \lr{Whisper Large-v3}،
یک مدل مبتنی بر \lr{Vosk}
و مدل
\lr{NVIDIA FastConformer Hybrid}
هستند.

هر یک از این مدل‌ها از رویکرد متفاوتی برای استخراج بازنمایی
صوت، مدل‌سازی توالی و رمزگشایی متن استفاده می‌کنند.
ازاین‌رو، بررسی معماری آن‌ها می‌تواند به درک بهتر تفاوت
عملکردشان در آزمایش‌های این پروژه کمک کند.


% ---------------------------------------------------------
\subsection{\lr{wav2vec 2.0} و \lr{XLSR}}

مدل
\lr{wav2vec 2.0}
یکی از رویکردهای مهم در زمینه
\textbf{یادگیری خودنظارتی گفتار}
\LTRfootnote{Self-Supervised \lr{Speech} Representation Learning}
است \cite{baevski2020wav2vec2}.

ایده اصلی این رویکرد آن است که بتوان از حجم بزرگی از داده‌های
صوتی فاقد متن برای یادگیری بازنمایی مناسب گفتار استفاده کرد
و سپس مدل حاصل را با مقدار محدودتری از داده‌های برچسب‌دار برای
وظیفه تشخیص گفتار تنظیم کرد.

در معماری \lr{wav2vec 2.0}، شکل موج خام صوت ابتدا توسط یک
رمزگذار ویژگی مبتنی بر شبکه‌های کانولوشنی پردازش می‌شود.
خروجی این بخش یک دنباله از بازنمایی‌های نهفته
\LTRfootnote{Latent Representations}
است.

بخشی از این بازنمایی‌ها در مرحله پیش‌آموزش به‌صورت تصادفی
پنهان می‌شوند.
سپس یک شبکه مبتنی بر
\lr{Transformer}
برای ایجاد بازنمایی‌های بسترآگاه از کل دنباله مورد استفاده
قرار می‌گیرد.

در کنار این مسیر، بازنمایی‌های نهفته توسط یک مؤلفه
\textbf{کوانتیزه‌سازی}
\LTRfootnote{Quantization}
به مجموعه‌ای از واحدهای گسسته تبدیل می‌شوند.
مدل در مرحله پیش‌آموزش تلاش می‌کند بازنمایی صحیح مربوط به
بخش پنهان‌شده را از میان چند گزینه نامرتبط تشخیص دهد.
این فرآیند با استفاده از یک
\textbf{تابع هدف تقابلی}
\LTRfootnote{Contrastive Objective}
انجام می‌شود.

پس از پیش‌آموزش، مدل می‌تواند با داده‌های گفتاری دارای
رونویسی برای وظیفه تشخیص گفتار تنظیم شود.
در این مرحله معمولاً از تابع زیان
\lr{CTC}
استفاده می‌شود:

\begin{center}
صوت خام
$\longrightarrow$
رمزگذار کانولوشنی
$\longrightarrow$
بازنمایی‌های نهفته
$\longrightarrow$
\lr{Transformer}
$\longrightarrow$
CTC
$\longrightarrow$
متن
\end{center}

یکی از مزایای مهم این روش آن است که بخش قابل توجهی از
بازنمایی گفتار پیش از استفاده از داده‌های دارای برچسب آموخته
می‌شود.
این ویژگی می‌تواند در زبان‌هایی که حجم داده رونویسی‌شده آن‌ها
نسبت به زبان‌های پرمنبع کمتر است، اهمیت بیشتری داشته باشد.

نسخه
\lr{XLSR}
\LTRfootnote{Cross-Lingual \lr{Speech} Representation}
این مفهوم را به محیط چندزبانه گسترش می‌دهد
\cite{conneau2020xlsr}.

در \lr{XLSR} یک مدل واحد با استفاده از صوت خام چندین زبان
پیش‌آموزش داده می‌شود.
هدف این است که مدل بتواند بازنمایی‌هایی را بیاموزد که بخشی از
ساختار آکوستیکی و واجی مشترک میان زبان‌ها را در خود جای دهند.
پس از این مرحله، مدل می‌تواند برای یک زبان مشخص با استفاده از
داده‌های برچسب‌دار همان زبان تنظیم شود.

مدل
\lr{wav2vec2-large-xlsr-persian-v3}
که در این پروژه مورد آزمایش قرار گرفته است، در همین خانواده
قرار دارد.
نتایج تجربی این مدل از نظر
\lr{WER}، \lr{CER}، \lr{RTF}
و مصرف منابع در فصل آزمایش‌ها با سایر مدل‌ها مقایسه خواهد شد.


% ---------------------------------------------------------
\subsection{\lr{Whisper}}

\lr{Whisper}
یک سامانه تشخیص و پردازش گفتار مبتنی بر معماری
\lr{Transformer}
است که با استفاده از حجم بزرگی از داده‌های صوتی همراه با
رونویسی آموزش داده شده است
\cite{radford2022whisper}.

بر خلاف مدل‌هایی مانند \lr{wav2vec 2.0} که بخش اصلی
پیش‌آموزش آن‌ها بر یادگیری خودنظارتی صوت متمرکز است،
\lr{Whisper}
در قالب یک مسئله نظارت‌شده چندوظیفه‌ای آموزش داده شده است.

ورودی مدل به‌صورت یک نمایش
\textbf{Log-Mel \lr{Spectrogram}}
از صوت آماده می‌شود.
معماری اصلی آن از دو بخش تشکیل شده است:

\begin{itemize}

    \item
    \textbf{رمزگذار}
    \LTRfootnote{Encoder}:
    ویژگی‌های صوت ورودی را پردازش کرده و بازنمایی سطح بالایی
    از سیگنال تولید می‌کند.

    \item
    \textbf{رمزگشا}
    \LTRfootnote{Decoder}:
    با استفاده از بازنمایی خروجی رمزگذار، متن را به‌صورت
    خودبازگشتی و توکن‌به‌توکن تولید می‌کند.

\end{itemize}

ساختار کلی مدل را می‌توان به شکل زیر نمایش داد:

\begin{center}
صوت
$\longrightarrow$
Log-Mel \lr{Spectrogram}
$\longrightarrow$
\lr{Transformer} Encoder
$\longrightarrow$
\lr{Transformer} Decoder
$\longrightarrow$
متن
\end{center}

یکی از ویژگی‌های مهم \lr{Whisper} آموزش چندوظیفه‌ای آن است.
مدل می‌تواند علاوه بر تشخیص گفتار، وظایفی مانند شناسایی زبان،
ترجمه گفتار به انگلیسی و پیش‌بینی اطلاعات زمانی را نیز
در قالب یک معماری مشترک انجام دهد.

این وظایف از طریق استفاده از توکن‌های کنترلی ویژه در ورودی
رمزگشا از یکدیگر تفکیک می‌شوند.

وجود رمزگشای خودبازگشتی به مدل امکان می‌دهد اطلاعات زبانی
توکن‌های تولیدشده قبلی را هنگام تولید توکن جدید در نظر بگیرد.
با این حال، از دیدگاه محاسباتی، تولید متوالی توکن‌ها می‌تواند
زمان استنتاج بیشتری نسبت به برخی معماری‌های غیرخودبازگشتی
ایجاد کند.

در این پروژه،
\lr{Whisper Large-v3}
به‌عنوان یکی از مدل‌های چندزبانه قدرتمند مورد ارزیابی قرار
گرفته است.
هدف از این ارزیابی بررسی این موضوع است که آیا توانایی عمومی
و چندزبانه این مدل در شرایط اجرای محلی فارسی، توازن مناسبی
میان دقت، سرعت و مصرف حافظه ایجاد می‌کند یا خیر.


% ---------------------------------------------------------
\subsection{\lr{Kaldi} و \lr{Vosk}}

\lr{Kaldi}
یک چارچوب متن‌باز برای پژوهش و توسعه سامانه‌های تشخیص گفتار
است \cite{povey2011kaldi}.
برخلاف معماری‌های کاملاً انتهابه‌انتهای جدید، چارچوب‌های
مبتنی بر \lr{Kaldi} امکان ترکیب چند مؤلفه مختلف برای
مدل‌سازی آکوستیکی، مدل‌سازی زبان و رمزگشایی را فراهم می‌کنند.

در یک ساختار کلاسیک مبتنی بر این رویکرد، مؤلفه‌هایی مانند
مدل آکوستیکی، واژه‌نامه تلفظ
\LTRfootnote{Pronunciation Lexicon}
و مدل زبانی در فرآیند رمزگشایی مشارکت دارند.

یکی از ابزارهای مهم در این نوع سامانه‌ها استفاده از
\textbf{مبدل‌های متناهی حالت وزن‌دار}
\LTRfootnote{Weighted Finite-State Transducers (WFST)}
است.
در این روش می‌توان اطلاعات مربوط به واحدهای آکوستیکی،
تلفظ واژه‌ها و مدل زبانی را در قالب گراف‌های قابل ترکیب
نمایش داد.

\lr{Vosk}
یک ابزار تشخیص گفتار با تمرکز بر استفاده عملی و اجرای
آفلاین است.
این ابزار مدل‌های ازپیش‌آموزش‌دیده برای زبان‌های مختلف ارائه
می‌دهد و قابلیت استفاده در سامانه‌های رومیزی و دستگاه‌های با
منابع محدود را دارد.

یکی از ویژگی‌های کاربردی \lr{Vosk} پشتیبانی از پردازش
جریانی
\LTRfootnote{Streaming}
است.
در این حالت، سامانه می‌تواند صوت را به‌صورت قطعات متوالی
دریافت کرده و بدون انتظار برای پایان یک فایل صوتی طولانی،
فرآیند تشخیص را انجام دهد.

اجرای برخی مدل‌های \lr{Vosk} بدون نیاز به پردازنده گرافیکی
نیز امکان‌پذیر است و این موضوع آن را برای بررسی سناریوهای
کم‌منبع مناسب می‌سازد.

در این پروژه مدل فارسی
\lr{vosk-model-fa-0.42}
به‌عنوان یکی از گزینه‌های قابل اجرای محلی مورد ارزیابی قرار
گرفته است.
این مدل به‌ویژه از منظر مصرف منابع و امکان اجرای مبتنی بر
\lr{CPU}
در مقایسه با مدل‌های عصبی بزرگ‌تر اهمیت دارد.


% ---------------------------------------------------------
\subsection{\lr{Conformer}}

معماری
\lr{Conformer}
با هدف ترکیب مزایای شبکه‌های کانولوشنی و
\lr{Transformer}
برای پردازش گفتار معرفی شده است
\cite{gulati2020conformer}.

در مدل‌های مبتنی بر \lr{Transformer}،
سازوکار توجه به مدل اجازه می‌دهد وابستگی میان بخش‌های دور
از یکدیگر در دنباله را مدل کند.
این ویژگی برای درک ساختار کلی گفتار مفید است.

از سوی دیگر، شبکه‌های کانولوشنی در استخراج الگوهای محلی
مانند تغییرات کوتاه‌مدت طیفی عملکرد مناسبی دارند.
معماری \lr{Conformer} تلاش می‌کند هر دو نوع اطلاعات را در
یک بلوک مشترک پردازش کند.

یک بلوک \lr{Conformer} به‌صورت مفهومی شامل اجزای زیر است:

\begin{center}
Feed-Forward
$\rightarrow$
Multi-Head Self-Attention
$\rightarrow$
Convolution
$\rightarrow$
Feed-Forward
\end{center}

به همراه اتصال‌های باقی‌مانده
\LTRfootnote{Residual Connections}
و نرمال‌سازی.

بخش توجه خودی
\LTRfootnote{Self-Attention}
به مدل کمک می‌کند اطلاعات وابسته به زمینه در بازه زمانی
طولانی‌تر را استخراج کند، در حالی که بخش کانولوشن برای
مدل‌سازی ساختارهای محلی گفتار مورد استفاده قرار می‌گیرد.

این ترکیب باعث شده است معماری \lr{Conformer} به یکی از
معماری‌های مهم در سامانه‌های مدرن تشخیص گفتار تبدیل شود.


% ---------------------------------------------------------
\subsection{\lr{FastConformer}}

با وجود عملکرد مناسب \lr{Conformer}، سازوکار توجه در دنباله‌های
طولانی می‌تواند هزینه محاسباتی و حافظه قابل توجهی ایجاد کند.
معماری
\lr{FastConformer}
با هدف کاهش بخشی از این هزینه طراحی شده است
\cite{rekesh2023fastconformer}.

یکی از تغییرات اصلی در \lr{FastConformer} استفاده از
نمونه‌برداری کاهشی بیشتر در ابتدای رمزگذار است.
کاهش طول دنباله باعث می‌شود لایه‌های بعدی توجه روی تعداد
کمتری از گام‌های زمانی اجرا شوند.

در این معماری از
\textbf{کانولوشن‌های عمقی تفکیک‌پذیر}
\LTRfootnote{Depthwise-Separable Convolutions}
نیز در بخش نمونه‌برداری کاهشی استفاده می‌شود.

به‌صورت مفهومی:

\begin{center}
ویژگی‌های صوتی
$\longrightarrow$
Downsampling
$\longrightarrow$
\lr{FastConformer} Encoder
$\longrightarrow$
Decoder
$\longrightarrow$
متن
\end{center}

هدف این تغییرات کاهش هزینه پردازش بدون از دست‌دادن بخش قابل
توجهی از دقت مدل است.

در پروژه حاضر، مدل فارسی

\begin{center}
\lr{nvidia/stt\_fa\_fastconformer\_hybrid\_large}
\end{center}

مورد ارزیابی قرار گرفته است.

این مدل از رمزگذار \lr{FastConformer} و یک ساختار
\textbf{Hybrid Transducer--CTC}
استفاده می‌کند.
به این معنا که مدل در مرحله آموزش با دو هدف
\lr{RNN-T} و \lr{CTC}
آموزش داده شده و امکان استفاده از سرهای رمزگشایی مختلف را
فراهم می‌کند.

این مدل به‌طور اختصاصی برای تشخیص گفتار فارسی آموزش داده شده
است و یکی از مدل‌های اصلی مورد بررسی در این پروژه به شمار
می‌رود.


% ---------------------------------------------------------
\subsection{\lr{RNN-T} و مدل‌های Hybrid}

\textbf{مبدل شبکه عصبی بازگشتی}
\LTRfootnote{Recurrent Neural Network Transducer (RNN-T)}
یک روش انتهابه‌انتها برای تبدیل دنباله ورودی به دنباله خروجی
است \cite{graves2012rnnt}.

بر خلاف \lr{CTC} که در شکل استاندارد خود فرض استقلال
شرطی بیشتری میان خروجی‌های زمانی دارد، \lr{RNN-T} می‌تواند
در تولید هر واحد خروجی، اطلاعات مربوط به خروجی‌های قبلی را
نیز در نظر بگیرد.

یک معماری \lr{RNN-T} معمولاً از سه بخش اصلی تشکیل می‌شود:

\begin{itemize}

    \item
    \textbf{Encoder}:
    سیگنال صوتی را به بازنمایی سطح بالاتر تبدیل می‌کند.

    \item
    \textbf{Prediction Network}:
    دنباله توکن‌های قبلی خروجی را پردازش می‌کند.

    \item
    \textbf{Joint Network}:
    اطلاعات رمزگذار و شبکه پیش‌بینی را ترکیب کرده و احتمال
    توکن بعدی را تولید می‌کند.

\end{itemize}

ساختار ساده‌شده آن را می‌توان به‌صورت زیر نمایش داد:

\begin{center}
\lr{Audio}
$\longrightarrow$
Encoder
$\longrightarrow$
Joint Network
$\longrightarrow$
Output Token
\end{center}

\begin{center}
Previous Tokens
$\longrightarrow$
Prediction Network
$\longrightarrow$
Joint Network
\end{center}

یکی از مزایای این ساختار امکان پردازش تدریجی دنباله و استفاده
در سناریوهای تشخیص گفتار جریانی است.

در مدل‌های
\textbf{Hybrid CTC--Transducer}
یک رمزگذار مشترک می‌تواند با بیش از یک تابع هدف آموزش داده شود.
به این ترتیب، مدل می‌تواند از مزایای آموزش مبتنی بر
\lr{CTC}
و \lr{Transducer}
به‌صورت هم‌زمان بهره گیرد.

مدل
\lr{stt\_fa\_fastconformer\_hybrid\_large}
که در این پژوهش استفاده شده است، نمونه‌ای از این ساختار
ترکیبی است.

وجود دو سر خروجی همچنین امکان انجام استنتاج با روش‌های
رمزگشایی مختلف را فراهم می‌کند.
این موضوع باید هنگام مقایسه نتایج تجربی مورد توجه قرار گیرد،
زیرا نوع سر رمزگشایی و تنظیمات آن می‌تواند بر دقت و سرعت
مدل اثرگذار باشد.


% ---------------------------------------------------------
\subsection{مقایسه مفهومی معماری‌های مورد بررسی}

مدل‌های بررسی‌شده در این پروژه از نظر روش آموزش،
نوع ورودی، معماری و فرآیند رمزگشایی تفاوت‌های قابل توجهی دارند.

\begin{table}[H]
\centering
\caption{مقایسه مفهومی خانواده مدل‌های تشخیص گفتار مورد بررسی}
\label{tab:asr_architecture_comparison}

\begin{tabular}{p{3.2cm} p{3.4cm} p{3.6cm} p{3.2cm}}
\toprule
\textbf{مدل} &
\textbf{رویکرد اصلی} &
\textbf{معماری} &
\textbf{روش خروجی} \\
\midrule

\lr{wav2vec2/XLSR} &
پیش‌آموزش خودنظارتی &
CNN + \lr{Transformer} &
CTC \\

\lr{Whisper} &
آموزش نظارت‌شده چندوظیفه‌ای &
\lr{Transformer} Encoder--Decoder &
Autoregressive \\

\lr{Vosk} &
سامانه تشخیص گفتار آفلاین &
مدل آکوستیکی + ساختار رمزگشایی &
Decoder-based \\

\lr{FastConformer Hybrid} &
تشخیص گفتار عصبی انتهابه‌انتها &
\lr{FastConformer} Encoder &
RNN-T / CTC \\

\bottomrule
\end{tabular}
\end{table}

این تفاوت‌های معماری سبب می‌شوند مقایسه مدل‌ها تنها بر اساس
یک معیار مانند نرخ خطای واژه کافی نباشد.
برای مثال، یک مدل ممکن است دقت بالاتری داشته باشد اما زمان
استنتاج یا مصرف حافظه آن برای اجرای هم‌زمان با مدل زبانی و
سامانه تبدیل متن به گفتار مناسب نباشد.

ازاین‌رو، در فصل آزمایش‌ها این معماری‌ها از چند جنبه شامل
دقت، سرعت، تأخیر و مصرف منابع با یکدیگر مقایسه خواهند شد.

% =========================================================
% 2-5
% =========================================================
\section{مدل‌های زبانی بزرگ}

\textbf{مدل‌های زبانی بزرگ}
\LTRfootnote{Large Language Models (LLMs)}
دسته‌ای از مدل‌های یادگیری عمیق هستند که با استفاده از حجم زیادی
از داده‌های متنی آموزش داده می‌شوند تا الگوهای آماری و ساختارهای
زبانی را بیاموزند.
این مدل‌ها می‌توانند وظایفی مانند تولید متن، پاسخ‌گویی به پرسش،
خلاصه‌سازی، ترجمه، استدلال متنی و گفت‌وگو را انجام دهند.

در یک چت‌بات صوتی، مدل زبانی نقش بخش مرکزی پردازش زبان را ایفا
می‌کند.
پس از آنکه مؤلفه تشخیص گفتار، صوت کاربر را به متن تبدیل کرد،
متن حاصل به مدل زبانی ارسال می‌شود.
مدل با در نظر گرفتن درخواست فعلی و در صورت نیاز تاریخچه مکالمه،
پاسخ مناسب را تولید می‌کند.
این پاسخ سپس به مؤلفه تبدیل متن به گفتار ارسال می‌شود.

به‌صورت ساده، نقش مدل زبانی در خط لوله پروژه حاضر را می‌توان
به شکل زیر نمایش داد:

\begin{center}
متن خروجی \lr{ASR}
$\longrightarrow$
مدل زبانی
$\longrightarrow$
پاسخ متنی
$\longrightarrow$
\lr{TTS}
\end{center}

با توجه به هدف اجرای محلی سامانه، انتخاب مدل زبانی تنها بر اساس
کیفیت پاسخ انجام نمی‌شود.
اندازه مدل، سرعت تولید توکن، زمان پاسخ و میزان حافظه مورد نیاز
نیز از معیارهای مهم در انتخاب مدل هستند.


% ---------------------------------------------------------
\subsection{معماری \lr{Transformer}}

بخش عمده مدل‌های زبانی بزرگ امروزی بر اساس معماری
\lr{Transformer}
ساخته شده‌اند.
این معماری با هدف پردازش مؤثر دنباله‌ها و مدل‌سازی وابستگی میان
اجزای مختلف آن‌ها معرفی شد \cite{vaswani2017attention}.

مؤلفه اصلی معماری \lr{Transformer} سازوکار
\textbf{توجه}
\LTRfootnote{Attention}
است.
ایده اصلی این سازوکار آن است که مدل هنگام پردازش یک توکن بتواند
اهمیت سایر توکن‌های موجود در دنباله را نیز در نظر بگیرد.

در سازوکار توجه، سه بازنمایی اصلی با نام‌های
\textbf{پرس‌وجو}
\LTRfootnote{Query}،
\textbf{کلید}
\LTRfootnote{Key}
و
\textbf{مقدار}
\LTRfootnote{Value}
برای هر توکن تولید می‌شوند.

توجه مقیاس‌شده ضرب داخلی
\LTRfootnote{Scaled Dot-Product Attention}
به‌صورت زیر تعریف می‌شود:

\begin{equation}
\mathrm{Attention}(Q,K,V)
=
\mathrm{softmax}
\left(
\frac{QK^T}{\sqrt{d_k}}
\right)V
\end{equation}

که در آن $Q$، $K$ و $V$ به‌ترتیب ماتریس‌های پرس‌وجو،
کلید و مقدار هستند و $d_k$ بعد بردارهای کلید را نشان می‌دهد.

عامل $\sqrt{d_k}$ برای کنترل مقیاس حاصل‌ضرب داخلی استفاده
می‌شود تا مقادیر ورودی تابع \lr{softmax} بیش از حد بزرگ نشوند.

در عمل، به‌جای یک سازوکار توجه واحد از
\textbf{توجه چندسری}
\LTRfootnote{Multi-Head Attention}
استفاده می‌شود.
در این روش، چندین سازوکار توجه به‌صورت موازی اجرا می‌شوند تا
مدل بتواند انواع مختلف روابط میان توکن‌ها را در فضاهای بازنمایی
متفاوت بیاموزد.

ساختار کلی یک بلوک \lr{Transformer} معمولاً شامل سازوکار توجه،
یک شبکه پیش‌خور
\LTRfootnote{Feed-Forward Network}
و اتصال‌های باقی‌مانده
\LTRfootnote{Residual Connections}
همراه با نرمال‌سازی است.


% ---------------------------------------------------------
\subsection{مدل‌های زبانی خودبازگشتی}

بسیاری از مدل‌های مورد استفاده در سامانه‌های گفت‌وگومحور،
از جمله مدل‌های بررسی‌شده در این پروژه، از نوع
\textbf{خودبازگشتی}
\LTRfootnote{Autoregressive}
هستند.

در چنین مدل‌هایی، متن به‌صورت توکن‌به‌توکن تولید می‌شود و
احتمال هر توکن جدید به توکن‌های قبلی وابسته است.

اگر دنباله توکن‌ها را با
$y_1,y_2,\ldots,y_T$
نمایش دهیم، احتمال دنباله کامل را می‌توان به‌صورت زیر نوشت:

\begin{equation}
P(y_1,\ldots,y_T)
=
\prod_{t=1}^{T}
P(y_t|y_1,\ldots,y_{t-1})
\end{equation}

در زمان تولید متن، مدل ابتدا ورودی کاربر و سایر اطلاعات موجود
در بافت را پردازش می‌کند و سپس احتمال توکن بعدی را محاسبه
می‌کند.
توکن انتخاب‌شده به دنباله اضافه شده و فرآیند تا رسیدن به
شرط پایان ادامه می‌یابد.

این نحوه تولید باعث می‌شود زمان پاسخ‌گویی مدل به عواملی مانند
طول ورودی، طول پاسخ و سرعت تولید هر توکن وابسته باشد.
به همین دلیل، در یک سامانه مکالمه‌ای معیار
\textbf{نرخ تولید توکن}
\LTRfootnote{Tokens per Second}
اهمیت زیادی دارد.


% ---------------------------------------------------------
\subsection{توکن‌سازی}

مدل‌های زبانی متن خام را مستقیماً به‌صورت رشته‌ای از حروف
پردازش نمی‌کنند.
پیش از ورود متن به مدل، متن توسط
\textbf{توکن‌ساز}
\LTRfootnote{Tokenizer}
به مجموعه‌ای از واحدهای عددی تبدیل می‌شود.

هر توکن ممکن است یک واژه کامل، بخشی از یک واژه، یک نویسه
یا حتی ترکیبی از چند نویسه باشد.
مدل در واقع روی شناسه‌های عددی این توکن‌ها عمل می‌کند.

اگر متن ورودی را با $X$ و تابع توکن‌سازی را با
$\mathcal{T}$ نمایش دهیم:

\begin{equation}
[t_1,t_2,\ldots,t_n]
=
\mathcal{T}(X)
\end{equation}

خواهد بود که $t_i$ شناسه توکن $i$ام است.

نوع توکن‌سازی می‌تواند بر تعداد توکن‌های لازم برای نمایش یک
زبان تأثیر بگذارد.
این موضوع برای زبان فارسی اهمیت دارد، زیرا شکل نوشتاری،
فاصله، نیم‌فاصله و ترکیب واژه‌ها می‌توانند باعث تفاوت در
تقسیم‌بندی متن شوند.

هرچه یک جمله به تعداد بیشتری توکن تقسیم شود، محاسبات بیشتری
برای پردازش و تولید آن لازم است.
ازاین‌رو، کارایی توکن‌ساز یکی از عوامل مؤثر بر هزینه استنتاج
مدل‌های زبانی است.


% ---------------------------------------------------------
\subsection{مدل‌های Decoder-Only}

یکی از رایج‌ترین معماری‌ها برای مدل‌های زبانی مولد،
معماری
\textbf{فقط رمزگشا}
\LTRfootnote{Decoder-Only}
است.

در این معماری، مدل از بلوک‌های متوالی \lr{Transformer} استفاده
می‌کند و برای تولید توکن بعدی تنها اجازه مشاهده توکن‌های قبلی
را دارد.
این محدودیت با استفاده از
\textbf{ماسک علّی}
\LTRfootnote{Causal Mask}
اعمال می‌شود.

در نتیجه، هنگام پیش‌بینی توکن $y_t$ مدل نمی‌تواند به توکن‌های
آینده دسترسی داشته باشد.

مدل‌های خانواده
\lr{Gemma 2}،
\lr{Qwen2.5}
و
\lr{Llama 3.1}
که مبنای مدل‌های مورد بررسی در این پروژه هستند، در دسته
مدل‌های مولد خودبازگشتی مبتنی بر \lr{Transformer} قرار می‌گیرند.
برای مثال، مستندات رسمی \lr{Llama} 3.1 آن را یک مدل زبانی
خودبازگشتی مبتنی بر معماری بهینه‌شده \lr{Transformer} معرفی
می‌کنند.


% ---------------------------------------------------------
\subsection{پیش‌آموزش و تنظیم دستورمحور}

فرآیند توسعه یک مدل زبانی معمولاً شامل چند مرحله است.
در مرحله
\textbf{پیش‌آموزش}
\LTRfootnote{Pre-training}
مدل با حجم بسیار بزرگی از داده متنی آموزش داده می‌شود تا
ساختار زبان، دانش عمومی و روابط آماری میان توکن‌ها را بیاموزد.

یکی از اهداف متداول در این مرحله پیش‌بینی توکن بعدی است.
به‌صورت ساده، مدل تلاش می‌کند تابع زیر را بیشینه کند:

\begin{equation}
\sum_{t=1}^{T}
\log
P(y_t|y_1,\ldots,y_{t-1})
\end{equation}

پس از پیش‌آموزش، مدل پایه ممکن است برای دنبال‌کردن دستورات
کاربر مناسب نباشد.
ازاین‌رو، مرحله دیگری تحت عنوان
\textbf{تنظیم دستورمحور}
\LTRfootnote{Instruction Tuning}
انجام می‌شود.

در این مرحله، نمونه‌هایی شامل دستور کاربر و پاسخ مناسب به مدل
ارائه می‌شوند تا رفتار آن برای کاربردهایی مانند پرسش و پاسخ و
گفت‌وگو سازگارتر شود.

مدل‌هایی که نام آن‌ها شامل عبارت
\lr{Instruct}
است، معمولاً نسخه‌هایی هستند که برای پیروی از دستور و تعامل
گفت‌وگویی تنظیم شده‌اند.
به همین دلیل، در پروژه حاضر نسخه‌های Instruct برای مقایسه
مدل‌های زبانی انتخاب شده‌اند.


% ---------------------------------------------------------
\subsection{پنجره بافت و مدیریت تاریخچه مکالمه}

یکی از قابلیت‌های مهم یک چت‌بات، توانایی استفاده از اطلاعات
مراحل قبلی مکالمه است.
برای مثال، در گفت‌وگوی زیر:

\begin{center}
کاربر: پایتخت فرانسه چیست؟ \\
سامانه: پاریس. \\
کاربر: جمعیت آن چقدر است؟
\end{center}

پرسش دوم تنها در صورتی به‌درستی تفسیر می‌شود که مدل بداند
کلمه «آن» به پاریس اشاره دارد.

برای حل این مسئله، بخشی از تاریخچه مکالمه همراه با پیام جدید
در ورودی مدل قرار می‌گیرد.

به‌صورت مفهومی می‌توان ورودی مدل در نوبت $t$ را چنین نمایش داد:

\begin{equation}
C_t
=
\{M_1,M_2,\ldots,M_{t-1},U_t\}
\end{equation}

که در آن $M_i$ پیام‌های پیشین و $U_t$ پیام فعلی کاربر است.

مدل‌های زبانی دارای محدودیتی در تعداد توکن‌هایی هستند که
می‌توانند در یک مرحله پردازش کنند.
این مقدار معمولاً
\textbf{طول بافت}
\LTRfootnote{Context Length}
نامیده می‌شود.

بنابراین در مکالمه‌های طولانی نمی‌توان بدون محدودیت تمام
تاریخچه را به مدل ارسال کرد.
راهکارهایی مانند حذف پیام‌های قدیمی، نگهداری تنها نوبت‌های
اخیر، خلاصه‌سازی تاریخچه یا ذخیره اطلاعات مهم می‌توانند برای
مدیریت این محدودیت استفاده شوند.

در سامانه این پروژه، مدیریت بافت مکالمه یکی از مؤلفه‌هایی است
که در مراحل تکمیلی پیاده‌سازی و ارزیابی خواهد شد.


% ---------------------------------------------------------
\subsection{فرآیند تولید پاسخ}

پس از پردازش ورودی، مدل برای هر گام یک توزیع احتمال روی
واژگان خود تولید می‌کند.

ساده‌ترین روش تولید، انتخاب توکنی با بیشترین احتمال است:

\begin{equation}
\hat{y}_t
=
\arg\max_y
P(y|y_{<t},X)
\end{equation}

با این حال، این روش ممکن است پاسخ‌های بیش از حد قطعی یا
تکراری ایجاد کند.
روش‌های نمونه‌برداری مختلفی برای کنترل تولید وجود دارند.

یکی از پارامترهای رایج
\textbf{دما}
\LTRfootnote{Temperature}
است.
با استفاده از دما، توزیع احتمال مدل پیش از نمونه‌برداری
تغییر داده می‌شود:

\begin{equation}
P_i'
=
\frac{
\exp(z_i/T)
}{
\sum_j \exp(z_j/T)
}
\end{equation}

که در آن $z_i$ مقدار لاجیت
\LTRfootnote{Logit}
توکن $i$ و $T$ مقدار دما است.

مقادیر کمتر دما توزیع را قطعی‌تر می‌کنند، در حالی که مقادیر
بالاتر احتمال انتخاب گزینه‌های متنوع‌تر را افزایش می‌دهند.

از روش‌های دیگری مانند
\lr{Top-k}
و
\lr{Top-p}
نیز می‌توان برای محدودکردن مجموعه توکن‌های قابل انتخاب
استفاده کرد.

برای ارزیابی منصفانه چند مدل، تنظیمات تولید باید تا حد امکان
ثابت نگه داشته شوند؛ زیرا تغییر پارامترهای تولید می‌تواند
کیفیت و سرعت پاسخ را تغییر دهد.


% ---------------------------------------------------------
\subsection{اجرای محلی مدل‌های زبانی}

اجرای محلی یک مدل زبانی با اجرای آن از طریق یک سرویس ابری
تفاوت مهمی دارد.
در اجرای محلی، وزن‌های مدل در حافظه سیستم یا حافظه گرافیکی
بارگذاری شده و عملیات استنتاج روی سخت‌افزار کاربر انجام می‌شود.

حافظه مورد نیاز یک مدل تا حد زیادی به تعداد پارامترها و
دقت عددی
\LTRfootnote{Numerical Precision}
استفاده‌شده برای ذخیره وزن‌ها وابسته است.

اگر مدلی دارای $N$ پارامتر باشد و هر پارامتر با $b$ بیت ذخیره
شود، حداقل حافظه مورد نیاز برای وزن‌ها را می‌توان تقریباً به شکل
زیر در نظر گرفت:

\begin{equation}
M
\approx
\frac{N \times b}{8}
\end{equation}

البته در عمل حافظه بیشتری برای مقادیر میانی، کش توجه و سایر
اجزای اجرای مدل مورد نیاز است.

به همین دلیل، یک مدل بزرگ‌تر لزوماً گزینه مناسب‌تری برای یک
چت‌بات محلی نیست.
مدل باید بتواند در کنار مدل تشخیص گفتار و مدل تبدیل متن به
گفتار روی سخت‌افزار موجود اجرا شود.


% ---------------------------------------------------------
\subsection{کوانتیزه‌سازی}

یکی از روش‌های متداول برای کاهش حافظه و هزینه اجرای مدل‌های
زبانی،
\textbf{کوانتیزه‌سازی}
\LTRfootnote{Quantization}
است.

در این روش، وزن‌ها یا فعال‌سازی‌های مدل با دقت عددی پایین‌تری
نمایش داده می‌شوند.
برای مثال، به‌جای استفاده از اعداد ممیز شناور 16 بیتی ممکن
است وزن‌ها با نمایش 8 بیتی یا 4 بیتی ذخیره شوند.

این روش می‌تواند مقدار حافظه مورد نیاز مدل را به میزان قابل
توجهی کاهش دهد و در برخی سخت‌افزارها سرعت اجرا را نیز افزایش
دهد.

با این حال، کاهش بیش از حد دقت عددی ممکن است کیفیت خروجی را
تحت تأثیر قرار دهد.
ازاین‌رو، انتخاب سطح کوانتیزه‌سازی نیز بخشی از توازن میان
کیفیت و هزینه محاسباتی است.


% ---------------------------------------------------------
\subsection{معیارهای کارایی در اجرای مدل زبانی}

در این پروژه، علاوه بر کیفیت پاسخ، عملکرد اجرایی مدل‌های زبانی
نیز مورد بررسی قرار گرفته است.

یکی از این معیارها
\textbf{زمان تأخیر}
است.
زمان تأخیر مشخص می‌کند تولید پاسخ برای یک ورودی چه مدت طول
می‌کشد.

معیار دیگر،
\textbf{نرخ تولید توکن}
است که به شکل تقریبی می‌توان آن را چنین تعریف کرد:

\begin{equation}
R_{\mathrm{token}}
=
\frac{N_{\mathrm{generated}}}
     {T_{\mathrm{generation}}}
\end{equation}

که در آن $N_{\mathrm{generated}}$ تعداد توکن‌های تولیدشده و
$T_{\mathrm{generation}}$ زمان تولید است.

مقدار بزرگ‌تر این معیار به معنی تولید سریع‌تر متن است.

میزان مصرف حافظه گرافیکی نیز در اجرای محلی اهمیت زیادی دارد.
در پروژه حاضر، مقدار بیشینه حافظه گرافیکی
\LTRfootnote{Peak \lr{GPU} Memory}
در کنار سرعت و کیفیت پاسخ برای مقایسه مدل‌ها استفاده شده است.


% ---------------------------------------------------------
\subsection{مدل‌های مورد بررسی در پروژه}

برای بخش مدل زبانی این پروژه، چند مدل با اندازه نسبتاً کوچک
یا متوسط انتخاب شده‌اند تا امکان اجرای آن‌ها روی سخت‌افزار شخصی
وجود داشته باشد.

مدل‌های مورد بررسی عبارت‌اند از:

\begin{itemize}

    \item
    \lr{gemma-2-2b-it}

    \item
    \lr{Qwen2.5-3B-Instruct}

    \item
    \lr{Qwen2.5-7B-Instruct}

    \item
    \lr{PartAI/\lr{Dorna2}-\lr{Llama}3.1-8B-Instruct}

\end{itemize}

خانواده \lr{Gemma 2} شامل مدل‌هایی با اندازه‌های مختلف است
و با هدف ارائه مدل‌های زبانی با اندازه عملی توسعه یافته است.
نسخه 2 میلیارد پارامتری این خانواده نیز ارائه شده و در طراحی
آن از روش‌هایی مانند ترکیب توجه محلی و سراسری و
\textbf{توجه گروهی پرس‌وجو}
\LTRfootnote{Grouped-Query Attention}
استفاده شده است \cite{gemma2report}.

خانواده \lr{Qwen2.5} شامل مدل‌هایی در اندازه‌های مختلف است و
نسخه‌های پایه و دستورمحور آن عرضه شده‌اند.
در گزارش فنی این خانواده، توسعه داده‌های پیش‌آموزش و
پس‌آموزش گسترده برای بهبود پیروی از دستور و توانایی‌های متنی
توضیح داده شده است \cite{qwen25report}.

مدل
\lr{Dorna2-Llama3.1-8B-Instruct}
بر پایه
\lr{Llama 3.1 8B Instruct}
ساخته شده و برای داده‌های فارسی آموزش یا تنظیم تکمیلی شده است.
این مدل از نوع Decoder-Only است و برای کاربردهای فارسی و
گفت‌وگویی ارائه شده است.

مدل پایه \lr{Llama 3.1 8B Instruct} نیز یک مدل خودبازگشتی
مبتنی بر \lr{Transformer} است که برای استفاده‌های گفت‌وگویی و
چندزبانه تنظیم شده است.

هدف از انتخاب این مدل‌ها مقایسه مدل‌هایی با اندازه و
ویژگی‌های متفاوت در شرایط اجرای محلی است.
نتایج عددی مربوط به کیفیت پاسخ، زمان تأخیر، سرعت تولید توکن
و مصرف حافظه در فصل آزمایش‌ها ارائه خواهند شد.


% =========================================================
% 2-6
% =========================================================
\section{تبدیل متن به گفتار}

\textbf{تبدیل متن به گفتار}
\LTRfootnote{\lr{Text}-to-\lr{Speech} (\lr{TTS})}
فرآیندی است که در آن یک سامانه رایانه‌ای متن ورودی را دریافت کرده
و یک سیگنال گفتاری قابل شنیدن تولید می‌کند.
هدف یک سامانه \lr{TTS} تنها تولید صوت نیست، بلکه خروجی باید
تا حد امکان طبیعی، روان و قابل فهم باشد.

در خط لوله پروژه حاضر، مؤلفه تبدیل متن به گفتار آخرین مرحله
پردازش است.
پس از آنکه مدل زبانی پاسخ متنی مناسب را تولید می‌کند،
این متن به مدل \lr{TTS} ارسال می‌شود تا پاسخ صوتی فارسی
برای کاربر ایجاد شود.

به‌صورت کلی، جایگاه \lr{TTS} در سامانه را می‌توان به شکل زیر
نمایش داد:

\begin{center}
پاسخ مدل زبانی
$\longrightarrow$
پردازش متن
$\longrightarrow$
مدل \lr{TTS}
$\longrightarrow$
سیگنال صوتی
$\longrightarrow$
پخش برای کاربر
\end{center}

در یک چت‌بات صوتی، کیفیت خروجی \lr{TTS} نقش مستقیمی در تجربه
کاربر دارد.
اگر گفتار تولیدشده دارای تلفظ نامناسب، مکث‌های غیرطبیعی،
ناپیوستگی یا نویز باشد، حتی یک پاسخ متنی صحیح ممکن است برای
کاربر دشوار یا نامطلوب باشد.
از سوی دیگر، مدل باید بتواند پاسخ را با سرعت کافی تولید کند
تا تأخیر کلی مکالمه افزایش زیادی پیدا نکند.


% ---------------------------------------------------------
\subsection{ساختار کلی سامانه‌های تبدیل متن به گفتار}

سامانه‌های تبدیل متن به گفتار را می‌توان به‌صورت مفهومی شامل
چند مرحله در نظر گرفت.

در مرحله نخست، متن ورودی پیش‌پردازش می‌شود.
این مرحله ممکن است شامل نرمال‌سازی متن، تبدیل اعداد و نشانه‌ها
به نمایش مناسب، مدیریت علائم نگارشی و در برخی سامانه‌ها تبدیل
متن به دنباله‌ای از واج‌ها
\LTRfootnote{Phonemes}
باشد.

سپس یک مدل صوتی
\LTRfootnote{Acoustic \lr{Model}}
نمایشی از ویژگی‌های گفتاری متن را تولید می‌کند.
در سامانه‌های کلاسیک‌تر، خروجی این بخش معمولاً یک طیف‌نگاشت
مانند Mel-\lr{Spectrogram} است.

در مرحله بعد، یک
\textbf{واکوُدر}
\LTRfootnote{\lr{Vocoder}}
این نمایش صوتی را به شکل موج نهایی تبدیل می‌کند.

بنابراین، ساختار کلاسیک را می‌توان به شکل زیر نمایش داد:

\begin{center}
\lr{Text}
$\longrightarrow$
\lr{Text} Processing
$\longrightarrow$
Acoustic \lr{Model}
$\longrightarrow$
\lr{Spectrogram}
$\longrightarrow$
\lr{Vocoder}
$\longrightarrow$
\lr{Waveform}
\end{center}

در نسل جدید سامانه‌های \lr{TTS}، بخشی یا تمام این مراحل می‌توانند
در یک معماری واحد ادغام شوند.
این رویکرد می‌تواند پیچیدگی خط لوله را کاهش داده و امکان آموزش
انتهابه‌انتهای سامانه را فراهم کند.


% ---------------------------------------------------------
\subsection{نرمال‌سازی متن}

متن تولیدشده توسط مدل زبانی ممکن است شامل اعداد، علائم اختصاری،
نشانه‌های نگارشی، نام‌های انگلیسی یا شکل‌های مختلف نوشتاری باشد.
سامانه \lr{TTS} باید بتواند این ورودی را به شکلی تبدیل کند که
تلفظ مناسب از آن تولید شود.

برای مثال، عدد

\begin{center}
1404
\end{center}

در متن گفتاری باید به شکلی معادل «هزار و چهارصد و چهار» تبدیل
شود و نه اینکه لزوماً هر رقم به‌صورت مستقل خوانده شود.

به همین دلیل،
\textbf{نرمال‌سازی متن}
\LTRfootnote{\lr{Text} Normalization}
یکی از مراحل مهم پیش از تولید گفتار است.

در زبان فارسی، مواردی مانند اعداد، تاریخ، علامت‌های اختصاری،
نیم‌فاصله و ترکیب واژه‌های فارسی و انگلیسی می‌توانند چالش‌های
خاصی برای این مرحله ایجاد کنند.


% ---------------------------------------------------------
\subsection{نمایش واجی}

برخی سامانه‌های تبدیل متن به گفتار، به‌جای استفاده مستقیم از
نویسه‌های متن، از نمایش واجی استفاده می‌کنند.

\textbf{واج}
\LTRfootnote{Phoneme}
کوچک‌ترین واحد صوتی تمایزدهنده در یک زبان است.
تبدیل متن به واج می‌تواند به مدل کمک کند تا تلفظ واژه‌ها را
مستقل‌تر از شکل نوشتاری آن‌ها بیاموزد.

فرآیند تبدیل نویسه یا متن به واج معمولاً با عنوان

\textbf{Grapheme-to-Phoneme}
\LTRfootnote{G2P}

شناخته می‌شود.

استفاده از نمایش واجی به‌ویژه زمانی مفید است که رابطه مستقیم
و یک‌به‌یک میان نوشتار و تلفظ وجود نداشته باشد.
در زبان فارسی نیز برخی واژه‌ها یا ترکیب‌های نوشتاری می‌توانند
با توجه به بافت، تلفظ متفاوتی داشته باشند.


% ---------------------------------------------------------
\subsection{معماری \lr{VITS}}

یکی از معماری‌های مهم تبدیل متن به گفتار،
\lr{VITS}
است \cite{kim2021vits}.

\lr{VITS}
یک معماری انتهابه‌انتها برای تولید گفتار است که اجزای مربوط به
مدل‌سازی صوت و تولید شکل موج را در یک چارچوب واحد ترکیب می‌کند.

در این معماری از
\textbf{استنتاج واریانسی}
\LTRfootnote{Variational Inference}،
\textbf{جریان‌های نرمال‌ساز}
\LTRfootnote{Normalizing Flows}
و آموزش
\textbf{خصمانه}
\LTRfootnote{Adversarial Training}
استفاده می‌شود.

یکی از ویژگی‌های مهم این معماری، مدل‌سازی رابطه چندبه‌یک میان
متن و گفتار است.
یک جمله یکسان می‌تواند با ریتم، آهنگ، سرعت و زیر و بمی متفاوت
بیان شود.
مدل \lr{VITS} تلاش می‌کند بخشی از این تنوع را در فضای نهفته
مدل‌سازی کند.

معماری را می‌توان به‌صورت مفهومی چنین نمایش داد:

\begin{center}
\lr{Text}
$\longrightarrow$
\lr{Text} Encoder
$\longrightarrow$
Latent Representation
$\longrightarrow$
Generator
$\longrightarrow$
\lr{Speech}
\end{center}

در طول آموزش، اطلاعات صوتی واقعی نیز توسط یک رمزگذار پسین
\LTRfootnote{Posterior Encoder}
پردازش می‌شوند و مدل از سازوکارهای واریانسی برای هم‌ترازکردن
نمایش متنی و صوتی استفاده می‌کند.

همچنین در \lr{VITS} از یک
\textbf{پیش‌بینی‌کننده مدت‌زمان تصادفی}
\LTRfootnote{Stochastic Duration Predictor}
استفاده می‌شود که به مدل امکان می‌دهد مدت زمان تلفظ واحدهای
زبانی را با تنوع بیشتری تولید کند.

این معماری به دلیل تولید مستقیم شکل موج، نیاز به یک واکودر
کاملاً مستقل در مرحله استنتاج را کاهش می‌دهد.


% ---------------------------------------------------------
\subsection{مدل \lr{MMS} برای تبدیل متن به گفتار}

پروژه
\lr{Massively Multilingual Speech}
یا
\lr{MMS}
با هدف توسعه فناوری‌های گفتاری برای تعداد بسیار زیادی از
زبان‌ها معرفی شده است \cite{pratap2023mms}.

این پروژه شامل مدل‌هایی برای تشخیص گفتار، شناسایی زبان و
تبدیل متن به گفتار در بیش از هزار زبان است.
این ویژگی آن را به یکی از منابع مهم برای زبان‌هایی تبدیل می‌کند
که در بسیاری از سامانه‌های تجاری یا مجموعه‌داده‌های بزرگ
پشتیبانی محدودتری دارند.

مدل فارسی مورد استفاده در این پروژه، که در آزمایش‌ها با عنوان

\begin{center}
\lr{VITS\_MMS\_FAS}
\end{center}

ذکر شده است، در این خانواده قرار دارد.

ساختار مدل‌های تولید گفتار \lr{MMS} بر مبنای معماری‌های عصبی
مانند \lr{VITS} است و هدف آن فراهم‌کردن امکان تولید گفتار
برای طیف بزرگی از زبان‌ها است.

در پروژه حاضر، این مدل به‌عنوان یکی از گزینه‌های قابل اجرای
محلی از نظر کیفیت صوت و سرعت استنتاج مورد بررسی قرار گرفته است.


% ---------------------------------------------------------
\subsection{سامانه \lr{Piper}}

\lr{Piper}
یک سامانه تبدیل متن به گفتار عصبی است که با تمرکز بر اجرای
سریع و محلی توسعه داده شده است.

مدل‌های صوتی \lr{Piper} بر پایه \lr{VITS} آموزش داده می‌شوند و
برای استنتاج به قالب
\lr{ONNX}
منتقل می‌شوند.
استفاده از \lr{ONNX Runtime} امکان اجرای مدل بر روی پلتفرم‌های
مختلف را فراهم می‌کند.

ساختار مفهومی اجرای \lr{Piper} را می‌توان چنین نمایش داد:

\begin{center}
\lr{Text}
$\longrightarrow$
\lr{Phonemization}
$\longrightarrow$
\lr{VITS}-based \lr{Model}
$\longrightarrow$
\lr{ONNX Runtime}
$\longrightarrow$
\lr{Speech}
\end{center}

تمرکز \lr{Piper} بر اجرای محلی باعث شده است که برای سامانه‌هایی
مانند دستیارهای صوتی و کاربردهای آفلاین گزینه مناسبی باشد.

مدل فارسی مورد استفاده در آزمایش‌های این پروژه با عنوان

\begin{center}
\lr{PIPER\_TTS\_FAS}
\end{center}

مورد ارزیابی قرار گرفته است.

از آنجا که سامانه نهایی پروژه نیز با هدف اجرای محلی طراحی شده
است، بررسی \lr{Piper} از نظر کیفیت و سرعت اهمیت ویژه‌ای دارد.


% ---------------------------------------------------------
\subsection{کیفیت ادراکی گفتار}

ارزیابی یک مدل \lr{TTS} تنها با معیارهای عددی مبتنی بر سیگنال
کافی نیست.
هدف نهایی این سامانه تولید گفتاری است که از دید شنونده طبیعی،
پیوسته و قابل فهم باشد.

یکی از معیارهای متداول در ارزیابی کیفیت صوت،
\textbf{میانگین امتیاز نظر}
\LTRfootnote{Mean Opinion Score (\lr{MOS})}
است.

در ارزیابی انسانی کلاسیک، چند شنونده کیفیت نمونه صوتی را روی
یک مقیاس عددی، معمولاً از 1 تا 5، امتیازدهی می‌کنند و میانگین
این امتیازها به‌عنوان \lr{MOS} گزارش می‌شود.

به‌صورت ساده:

\begin{equation}
MOS
=
\frac{1}{N}
\sum_{i=1}^{N} r_i
\end{equation}

که در آن $r_i$ امتیاز شنونده $i$ام و $N$ تعداد امتیازها است.

مقدار بیشتر \lr{MOS} نشان‌دهنده کیفیت ادراکی بالاتر است.


% ---------------------------------------------------------
\subsection{ارزیابی خودکار با \lr{NISQA}}

در این پروژه برای تخمین خودکار کیفیت گفتار تولیدشده از مدل
\lr{NISQA}
استفاده شده است.

\lr{NISQA}
یک مدل یادگیری عمیق برای پیش‌بینی کیفیت گفتار بدون نیاز به
سیگنال مرجع است \cite{mittag2021nisqa}.

این مدل علاوه بر تخمین امتیاز کلی کیفیت، چند بعد دیگر از
کیفیت صوت را نیز پیش‌بینی می‌کند:

\begin{itemize}

    \item
    \textbf{\lr{Noise}}:
    میزان اختلال ناشی از نویز موجود در صوت.

    \item
    \textbf{\lr{Discontinuity}}:
    میزان ناپیوستگی یا قطع‌شدگی در گفتار.

    \item
    \textbf{\lr{Coloration}}:
    تغییرات یا اعوجاج در ویژگی‌های طیفی و رنگ صوت.

    \item
    \textbf{\lr{Loudness}}:
    کیفیت مرتبط با سطح بلندی صوت.

\end{itemize}

این معیارها امکان بررسی جزئی‌تر علت کاهش کیفیت یک خروجی صوتی
را فراهم می‌کنند.

در آزمایش‌های این پروژه، نمونه‌های تولیدشده توسط هر مدل
\lr{TTS} با استفاده از \lr{NISQA} ارزیابی شده‌اند و امتیاز
کلی و ابعاد مختلف کیفیت برای مقایسه مدل‌ها مورد استفاده قرار
گرفته‌اند.


% ---------------------------------------------------------
\subsection{ارزیابی قابل‌فهم‌بودن گفتار با \lr{ASR}}

یک مدل \lr{TTS} ممکن است از نظر شنیداری طبیعی به نظر برسد،
اما واژه‌ها را به شکل نادرست یا مبهم تولید کند.
بنابراین علاوه بر کیفیت ادراکی، میزان
\textbf{قابل‌فهم‌بودن گفتار}
\LTRfootnote{\lr{Speech} Intelligibility}
نیز اهمیت دارد.

یکی از روش‌های خودکار برای تخمین قابل‌فهم‌بودن خروجی این است
که صوت تولیدشده مجدداً توسط یک سامانه تشخیص گفتار به متن
تبدیل شود.

فرآیند ارزیابی به‌صورت زیر است:

\begin{center}
متن اصلی
$\longrightarrow$
\lr{TTS}
$\longrightarrow$
صوت تولیدشده
$\longrightarrow$
\lr{ASR}
$\longrightarrow$
متن بازشناسی‌شده
\end{center}

سپس متن بازشناسی‌شده با متن اولیه مقایسه می‌شود و معیارهایی
مانند \lr{WER} و \lr{CER} محاسبه می‌شوند.

اگر خروجی \lr{TTS} واضح و قابل فهم باشد، انتظار می‌رود مدل
تشخیص گفتار بتواند بخش بیشتری از متن اولیه را بازسازی کند و
در نتیجه مقادیر \lr{WER} و \lr{CER} کمتر باشند.

با این حال، این روش دارای یک محدودیت مهم است:
نتیجه ارزیابی نه‌تنها به کیفیت \lr{TTS} بلکه به دقت مدل
\lr{ASR} مورد استفاده نیز وابسته است.
بنابراین خطای مدل تشخیص گفتار می‌تواند مستقیماً بر مقدار نهایی
\lr{WER} یا \lr{CER} اثر بگذارد.

این محدودیت در تفسیر نتایج آزمایش‌های \lr{TTS} این پروژه
در نظر گرفته خواهد شد.


% ---------------------------------------------------------
\subsection{ضریب زمان واقعی در \lr{TTS}}

مانند مدل تشخیص گفتار، برای ارزیابی سرعت مدل تبدیل متن به
گفتار نیز می‌توان از
\textbf{ضریب زمان واقعی}
استفاده کرد.

در این حالت:

\begin{equation}
RTF_{\\mathrm{TTS}}
=
\frac{T_{\mathrm{synthesis}}}
     {T_{\mathrm{generated\ audio}}}
\end{equation}

که در آن $T_{\mathrm{synthesis}}$ زمان مورد نیاز برای تولید
صوت و $T_{\mathrm{generated\ audio}}$ مدت زمان صوت تولیدشده
است.

اگر مقدار این معیار کمتر از 1 باشد، مدل می‌تواند گفتار را
سریع‌تر از زمان واقعی تولید کند.

برای مثال اگر تولید یک صوت پنج‌ثانیه‌ای تنها نیم ثانیه زمان
ببرد:

\begin{equation}
RTF_{\\mathrm{TTS}}
=
\frac{0.5}{5}
=
0.1
\end{equation}

خواهد بود.


% ---------------------------------------------------------
\subsection{ضریب افزایش سرعت}

یکی دیگر از روش‌های بیان سرعت مدل، استفاده از
\textbf{\lr{Speedup} نسبت به زمان واقعی}
است.

این معیار را می‌توان از معکوس \lr{RTF} محاسبه کرد:

\begin{equation}
Speedup
=
\frac{1}{RTF}
\end{equation}

برای مثال، اگر:

\begin{equation}
RTF = 0.05
\end{equation}

باشد، سرعت مدل نسبت به زمان واقعی برابر است با:

\begin{equation}
Speedup
=
20\times
\end{equation}

یعنی مدل می‌تواند صوت را حدود بیست برابر سریع‌تر از طول واقعی
آن تولید کند.


% ---------------------------------------------------------
\subsection{زمان تأخیر در تولید گفتار}

در یک سامانه مکالمه‌ای، علاوه بر \lr{RTF}، مدت زمان انتظار
کاربر برای آغاز یا تکمیل تولید پاسخ نیز اهمیت دارد.

اگر زمان تولید نمونه $i$ام را با $t_i$ و تعداد نمونه‌ها را با
$N$ نمایش دهیم، متوسط زمان تأخیر به شکل زیر است:

\begin{equation}
\mathrm{Average\ Latency}
=
\frac{1}{N}
\sum_{i=1}^{N} t_i
\end{equation}

یک مدل ممکن است \lr{RTF} بسیار مناسبی داشته باشد، اما برای
نمونه‌های کوتاه زمان شروع یا آماده‌سازی بیشتری نیاز داشته باشد.
ازاین‌رو استفاده هم‌زمان از \lr{RTF} و زمان تأخیر تصویر دقیق‌تری
از عملکرد اجرایی مدل ارائه می‌دهد.


% ---------------------------------------------------------
\subsection{مدل‌های \lr{TTS} مورد بررسی در پروژه}

در پروژه حاضر سه سامانه اصلی تبدیل متن به گفتار فارسی مورد
ارزیابی قرار گرفته‌اند:

\begin{itemize}

    \item
    \lr{Kamtera\_VITS}

    \item
    \lr{VITS\_MMS\_FAS}

    \item
    \lr{PIPER\_TTS\_FAS}

\end{itemize}

مقایسه این مدل‌ها بر اساس دو گروه اصلی از معیارها انجام شده
است.

گروه اول مربوط به کیفیت خروجی است و شامل امتیاز تخمینی
\lr{MOS} و ابعاد کیفیت پیش‌بینی‌شده توسط \lr{NISQA} می‌شود.

گروه دوم به قابلیت استفاده عملی از مدل مربوط است و شامل
\lr{WER}، \lr{CER}، زمان تأخیر و ضریب زمان واقعی است.

این رویکرد امکان بررسی هم‌زمان دو نیاز اصلی سامانه را فراهم
می‌کند:
کیفیت مناسب گفتار و سرعت کافی برای استفاده در مکالمه تعاملی.

نتایج عددی مدل‌های فوق در فصل آزمایش‌ها ارائه می‌شوند و
دلایل انتخاب مدل نهایی در فصل بحث و تحلیل بررسی خواهند شد.

% =========================================================
% 2-7
% =========================================================
\section{تشخیص فعالیت صوتی و پردازش نویز}

در یک سامانه گفت‌وگوی صوتی، ورودی میکروفون معمولاً شامل
ترکیبی از گفتار، سکوت و نویز محیط است.
اگر کل جریان صوتی بدون هیچ‌گونه جداسازی در اختیار مدل تشخیص
گفتار قرار گیرد، ممکن است بخش‌های طولانی سکوت، صداهای محیطی
یا نویزهای غیرگفتاری نیز پردازش شوند.
این موضوع می‌تواند علاوه بر افزایش هزینه محاسباتی، باعث کاهش
دقت تشخیص و افزایش تأخیر سامانه شود.

به همین دلیل، پیش از اجرای بخش
\lr{ASR}
معمولاً از پردازش‌هایی مانند
\textbf{تشخیص فعالیت صوتی}
\LTRfootnote{Voice Activity Detection (VAD)}
و در صورت نیاز کاهش نویز استفاده می‌شود.

جایگاه این پردازش‌ها در خط لوله را می‌توان به شکل زیر نمایش داد:

\begin{center}
Microphone
$\longrightarrow$
VAD
$\longrightarrow$
\lr{Noise} Processing
$\longrightarrow$
\lr{ASR}
$\longrightarrow$
\lr{Text}
\end{center}

هدف این بخش معرفی مفاهیم اصلی مربوط به تشخیص فعالیت صوتی،
تشخیص نقطه پایان گفتار، انواع نویز و روش‌های کلی کاهش اثر
نویز بر عملکرد سامانه است.


% ---------------------------------------------------------
\subsection{تشخیص فعالیت صوتی}

وظیفه اصلی یک سامانه
\lr{VAD}
تعیین این است که آیا در یک بازه زمانی مشخص از سیگنال صوتی،
گفتار انسانی وجود دارد یا خیر.

اگر تصمیم سامانه را در زمان $t$ با $V(t)$ نمایش دهیم:

\begin{equation}
V(t)
=
\begin{cases}
1 & \mbox{\rl{در صورت وجود گفتار}} \\
0 & \mbox{\rl{در صورت نبود گفتار}}
\end{cases}
\end{equation}

خواهد بود.

در ساده‌ترین حالت، یک الگوریتم VAD می‌تواند از ویژگی‌هایی
مانند انرژی سیگنال برای تشخیص گفتار استفاده کند.
برای مثال اگر انرژی یک قاب صوتی از مقدار آستانه مشخصی بیشتر
باشد، آن قاب می‌تواند به‌عنوان گفتار در نظر گرفته شود.

اگر انرژی قاب $i$ام را با $E_i$ و مقدار آستانه را با
$\theta$ نمایش دهیم، یک تصمیم ساده می‌تواند به شکل زیر باشد:

\begin{equation}
V_i
=
\begin{cases}
1 & E_i > \theta \\
0 & E_i \leq \theta
\end{cases}
\end{equation}

با این حال، چنین رویکرد ساده‌ای در محیط‌های نویزی عملکرد
محدودی دارد؛ زیرا نویز محیط نیز می‌تواند انرژی بالایی داشته
باشد.
به همین دلیل، سامانه‌های جدید VAD معمولاً از مجموعه‌ای از
ویژگی‌های صوتی یا مدل‌های یادگیری ماشین و شبکه‌های عصبی
استفاده می‌کنند.


% ---------------------------------------------------------
\subsection{تشخیص شروع و پایان گفتار}

در یک چت‌بات صوتی، صرفاً دانستن این که در یک قاب گفتار وجود
دارد کافی نیست.
سامانه باید بتواند تشخیص دهد که کاربر چه زمانی صحبت‌کردن را
آغاز کرده و چه زمانی جمله خود را به پایان رسانده است.

این مسئله معمولاً با مفهوم
\textbf{تشخیص نقطه پایان}
\LTRfootnote{Endpoint Detection}
مرتبط است.

یک پیاده‌سازی معمولی ممکن است از چند پارامتر استفاده کند:

\begin{itemize}

    \item حداقل مدت زمان گفتار برای شروع ثبت جمله؛

    \item حداقل مدت زمان سکوت برای تشخیص پایان جمله؛

    \item حداکثر طول مجاز یک بخش گفتاری؛

    \item مقدار حاشیه زمانی قبل و بعد از گفتار.

\end{itemize}

برای مثال، اگر کاربر برای مدت کوتاهی در وسط جمله مکث کند،
سامانه نباید فوراً فرض کند که جمله پایان یافته است.
در مقابل، اگر مدت سکوت بیش از حد طولانی در نظر گرفته شود،
کاربر پس از پایان صحبت مجبور خواهد بود زمان بیشتری منتظر
شروع پردازش بماند.

بنابراین انتخاب مدت زمان سکوت برای تشخیص پایان گفتار مستقیماً
بر تجربه کاربری و تأخیر انتهابه‌انتهای سامانه اثر می‌گذارد.


% ---------------------------------------------------------
\subsection{خطاهای تشخیص فعالیت صوتی}

عملکرد یک سامانه VAD را می‌توان از چند نوع خطا بررسی کرد.

یکی از خطاها زمانی رخ می‌دهد که سامانه بخش غیرگفتاری را
به‌عنوان گفتار تشخیص دهد.
این حالت را می‌توان
\textbf{تشخیص مثبت کاذب}
\LTRfootnote{False Positive}
در نظر گرفت.

در این حالت، نویز یا سکوت ممکن است به مدل تشخیص گفتار ارسال
شود و موجب پردازش غیرضروری یا حتی تولید متن اشتباه شود.

حالت دیگر زمانی است که بخشی از گفتار واقعی به اشتباه
غیرگفتاری تشخیص داده شود.
این وضعیت
\textbf{منفی کاذب}
\LTRfootnote{False Negative}
نام دارد.

این خطا می‌تواند شدیدتر باشد، زیرا ممکن است بخشی از جمله
کاربر حذف شده و در نتیجه متن تولیدشده توسط \lr{ASR} ناقص باشد.

اگر تعداد نمونه‌های گفتاری صحیح تشخیص‌داده‌شده را با $TP$،
نمونه‌های غیرگفتاری که اشتباه گفتار تشخیص داده شده‌اند را با
$FP$ و گفتارهای ازدست‌رفته را با $FN$ نمایش دهیم، معیارهایی
مانند دقت و بازخوانی را می‌توان به‌صورت زیر تعریف کرد:

\begin{equation}
Precision
=
\frac{TP}{TP+FP}
\end{equation}

\begin{equation}
Recall
=
\frac{TP}{TP+FN}
\end{equation}

و امتیاز $F_1$ برابر است با:

\begin{equation}
F_1
=
2
\frac{Precision \times Recall}
     {Precision + Recall}
\end{equation}

این معیارها در صورت طراحی یک مجموعه آزمایشی مناسب می‌توانند
برای ارزیابی VAD مورد استفاده قرار گیرند.


% ---------------------------------------------------------
\subsection{نقش VAD در تأخیر سامانه}

تشخیص فعالیت صوتی تنها برای حذف سکوت‌ها استفاده نمی‌شود،
بلکه نقش مهمی در زمان پاسخ سامانه نیز دارد.

فرض شود کاربر جمله خود را در زمان $t_e$ به پایان می‌رساند،
اما VAD پایان گفتار را در زمان $t_d$ تشخیص می‌دهد.
در این صورت، تأخیر ناشی از تشخیص پایان گفتار برابر است با:

\begin{equation}
T_{\mathrm{endpoint}}
=
t_d - t_e
\end{equation}

این مقدار بخشی از تأخیر کلی سامانه خواهد بود.

بنابراین:

\begin{equation}
T_{\mathrm{total}}
=
T_{\mathrm{endpoint}}
+
T_{\\mathrm{ASR}}
+
T_{\\mathrm{LLM}}
+
T_{\\mathrm{TTS}}
+
T_{\mathrm{other}}
\end{equation}

در نتیجه، حتی اگر مدل‌های \lr{ASR}، \lr{LLM} و \lr{TTS} سرعت بالایی داشته
باشند، تنظیم نامناسب VAD می‌تواند تأخیر قابل توجهی به تجربه
مکالمه اضافه کند.


% ---------------------------------------------------------
\subsection{نویز در سیگنال گفتار}

نویز به هر مؤلفه ناخواسته در سیگنال صوتی گفته می‌شود که
می‌تواند کیفیت گفتار یا عملکرد پردازش‌های بعدی را کاهش دهد.

نویزها را می‌توان از دیدگاه‌های مختلف دسته‌بندی کرد.

یکی از انواع رایج،
\textbf{نویز ایستا}
\LTRfootnote{Stationary \lr{Noise}}
است.
ویژگی‌های آماری این نوع نویز در طول زمان تغییرات محدودی دارد.
صدای ثابت فن رایانه یا سیستم تهویه نمونه‌ای از این حالت است.

در مقابل،
\textbf{نویز غیرایستا}
\LTRfootnote{Non-stationary \lr{Noise}}
ویژگی‌هایی دارد که در طول زمان تغییر می‌کنند.
صدای صحبت سایر افراد، صدای خودرو یا صدای موسیقی نمونه‌هایی از
این نوع نویز هستند.

به‌طور کلی نویز غیرایستا چالش بیشتری برای سامانه‌های حذف نویز
و تشخیص گفتار ایجاد می‌کند.


% ---------------------------------------------------------
\subsection{نسبت سیگنال به نویز}

همان‌طور که پیش‌تر معرفی شد، یکی از معیارهای متداول برای
اندازه‌گیری شدت نسبی گفتار و نویز،
\textbf{نسبت سیگنال به نویز}
\LTRfootnote{Signal-to-\lr{Noise} Ratio (SNR)}
است:

\begin{equation}
\mathrm{SNR}_{dB}
=
10
\log_{10}
\left(
\frac{P_{\mathrm{speech}}}
     {P_{\mathrm{noise}}}
\right)
\end{equation}

مقدار بالاتر SNR معمولاً به معنی وجود سیگنال گفتاری واضح‌تر
نسبت به نویز است.

برای ارزیابی مقاومت یک مدل \lr{ASR} در برابر نویز، می‌توان یک
سیگنال گفتاری تمیز را با نویز در سطوح مختلف SNR ترکیب کرد
و سپس تغییر معیارهایی مانند \lr{WER} و \lr{CER} را اندازه‌گیری نمود.

برای مثال، یک آزمایش می‌تواند شامل شرایط زیر باشد:

\begin{itemize}

    \item صوت تمیز؛

    \item \lr{SNR = 20 dB}؛

    \item \lr{SNR = 10 dB}؛

    \item \lr{SNR = 5 dB}؛

    \item \lr{SNR = 0 dB}.

\end{itemize}

این نوع آزمایش امکان بررسی میزان کاهش دقت مدل با افزایش شدت
نویز را فراهم می‌کند.


% ---------------------------------------------------------
\subsection{کاهش نویز}

\textbf{کاهش نویز}
\LTRfootnote{\lr{Noise} Reduction}
به مجموعه روش‌هایی گفته می‌شود که با هدف کاهش مؤلفه‌های
ناخواسته صوت پیش از پردازش‌های بعدی اجرا می‌شوند.

روش‌های کاهش نویز می‌توانند از تکنیک‌های کلاسیک پردازش سیگنال
تا مدل‌های عمیق یادگیری ماشین گسترده باشند.

یکی از روش‌های کلاسیک،
\textbf{تفریق طیفی}
\LTRfootnote{Spectral Subtraction}
است.
در این روش، طیف نویز تخمین زده شده و از طیف سیگنال آلوده
کم می‌شود.

به‌صورت ساده:

\begin{equation}
|\hat{S}(f)|
=
|Y(f)| - |\hat{N}(f)|
\end{equation}

که در آن $Y(f)$ طیف سیگنال نویزی،
$\hat{N}(f)$ تخمین طیف نویز و
$\hat{S}(f)$ تخمین طیف گفتار تمیز است.

روش دیگر استفاده از
\textbf{فیلتر وینر}
\LTRfootnote{Wiener Filter}
است که تلاش می‌کند خطای میان سیگنال تخمینی و سیگنال اصلی را
به حداقل برساند.

در روش‌های جدیدتر، شبکه‌های عصبی برای تخمین ماسک طیفی یا
مستقیماً برای بازسازی سیگنال گفتاری تمیز استفاده می‌شوند.


% ---------------------------------------------------------
\subsection{کاهش نویز لزوماً باعث بهبود \lr{ASR} نمی‌شود}

هرچند هدف از کاهش نویز ایجاد صوت واضح‌تر است، اعمال چنین
پردازشی همیشه منجر به بهبود تشخیص گفتار نمی‌شود.

یک الگوریتم کاهش نویز ممکن است علاوه بر نویز، بخشی از مؤلفه‌های
مفید گفتار را نیز حذف یا تغییر دهد.
این پدیده می‌تواند نوع جدیدی از اعوجاج ایجاد کند که برای مدل
\lr{ASR} ناآشنا باشد.

از سوی دیگر، بسیاری از مدل‌های جدید تشخیص گفتار با داده‌های
متنوع و تا حدی نویزی آموزش داده شده‌اند و ممکن است بدون
پیش‌پردازش شدید نیز مقاومت مناسبی داشته باشند.

بنابراین اثر کاهش نویز باید به‌صورت تجربی بررسی شود.
یک روش مناسب می‌تواند مقایسه \lr{WER} در دو حالت باشد:

\begin{center}
\lr{Noise} + \lr{ASR}
\end{center}

و:

\begin{center}
\lr{Noise} + Denoising + \lr{ASR}
\end{center}

اگر مقدار \lr{WER} در حالت دوم کمتر باشد، پردازش نویز برای آن
سناریوی خاص مفید بوده است.
در غیر این صورت، ممکن است هزینه و تأخیر اضافه‌شده توسط
ماژول کاهش نویز ارزش استفاده نداشته باشد.


% ---------------------------------------------------------
\subsection{پردازش بلادرنگ}

در یک چت‌بات صوتی، VAD و کاهش نویز باید تا حد امکان با
پردازش بلادرنگ سازگار باشند.

اگر این مؤلفه‌ها به زمان قابل توجهی برای پردازش نیاز داشته
باشند، تأخیر کلی سامانه افزایش خواهد یافت.

اگر مدت زمان لازم برای پردازش یک قطعه صوتی برابر با
$T_p$ و مدت زمان خود صوت برابر با $T_a$ باشد، می‌توان مانند
سایر ماژول‌های صوتی ضریب زمان واقعی را تعریف کرد:

\begin{equation}
RTF_{\mathrm{preprocess}}
=
\frac{T_p}{T_a}
\end{equation}

برای استفاده تعاملی، مطلوب است که این مقدار تا حد امکان کمتر
از 1 باشد.


% ---------------------------------------------------------
\subsection{جایگاه این مؤلفه‌ها در پروژه}

در نسخه پایه سامانه این پروژه، خط لوله اصلی
\lr{ASR--\lr{LLM}--\lr{TTS}}
به‌صورت محلی پیاده‌سازی شده است.

تشخیص فعالیت صوتی و پردازش نویز به‌عنوان مؤلفه‌های تکمیلی
در نظر گرفته شده‌اند تا در ادامه پروژه به خط لوله اضافه شوند.

معماری مورد نظر برای این بخش به‌صورت مفهومی شامل مراحل زیر
است:

\begin{center}
Microphone
$\longrightarrow$
VAD
$\longrightarrow$
Optional Denoising
$\longrightarrow$
\lr{ASR}
$\longrightarrow$
\lr{Text} Processing
\end{center}

پس از پیاده‌سازی، عملکرد این مؤلفه‌ها باید نه‌تنها به‌صورت
مستقل، بلکه از نظر تأثیر آن‌ها بر \lr{WER}، تأخیر و کیفیت
انتهابه‌انتهای سامانه نیز بررسی شود.

به‌ویژه سه پرسش عملی اهمیت دارند:

\begin{itemize}

    \item آیا VAD می‌تواند شروع و پایان گفتار را بدون حذف
    بخش‌های مهم جمله به‌درستی تشخیص دهد؟

    \item آیا کاهش نویز واقعاً دقت \lr{ASR} را در شرایط نویزی
    بهبود می‌دهد؟

    \item هزینه زمانی این پردازش‌ها چه اثری بر تأخیر کلی
    مکالمه دارد؟

\end{itemize}

پاسخ تجربی به این پرسش‌ها در فصل آزمایش‌ها ارائه خواهد شد.
% =========================================================
% 2-8
% =========================================================
\section{معماری سامانه‌های گفت‌وگوی صوتی}

یک سامانه گفت‌وگوی صوتی از چند مؤلفه متوالی تشکیل می‌شود که
هر کدام وظیفه مشخصی در تبدیل گفتار کاربر به پاسخ صوتی دارند.
در ساده‌ترین حالت، جریان پردازش شامل دریافت صوت، تشخیص گفتار،
پردازش زبان، تولید پاسخ و در نهایت تبدیل متن به گفتار است.

در پروژه حاضر، معماری کلی سامانه را می‌توان به‌صورت زیر در نظر گرفت:

\begin{center}
Microphone
$\longrightarrow$
VAD / \lr{Audio} Preprocessing
$\longrightarrow$
\lr{ASR}
$\longrightarrow$
\lr{Text} Processing
$\longrightarrow$
Conversation Context
$\longrightarrow$
\lr{LLM}
$\longrightarrow$
\lr{TTS}
$\longrightarrow$
Speaker
\end{center}

هر یک از این مراحل می‌تواند بر عملکرد کلی سامانه اثر بگذارد.
به‌عنوان مثال، اگر خروجی \lr{ASR} دارای خطای زیاد باشد، متن
نامناسبی در اختیار مدل زبانی قرار خواهد گرفت.
در نتیجه، مدل زبانی ممکن است پاسخ نادرستی تولید کند، حتی اگر
خود مدل در شرایط متنی صحیح عملکرد مناسبی داشته باشد.

به همین ترتیب، اگر پاسخ مدل زبانی صحیح باشد اما مدل \lr{TTS}
تلفظ مناسبی ارائه ندهد، تجربه کاربر در مرحله نهایی تحت تأثیر
قرار خواهد گرفت.

ازاین‌رو، یک سامانه گفت‌وگوی صوتی باید به‌صورت یک زنجیره
وابسته به هم مورد بررسی قرار گیرد و ارزیابی تنها یک جزء برای
قضاوت درباره کیفیت کل سامانه کافی نیست.


% ---------------------------------------------------------
\subsection{پردازش ترتیبی}

یکی از ساده‌ترین روش‌های طراحی چنین سامانه‌ای،
\textbf{پردازش ترتیبی}
\LTRfootnote{Sequential Processing}
است.

در این معماری، هر مرحله تنها پس از پایان کامل مرحله قبلی آغاز
می‌شود.

برای مثال:

\begin{enumerate}
    \item کاربر صحبت می‌کند.
    \item صوت کامل در اختیار مدل \lr{ASR} قرار می‌گیرد.
    \item پس از پایان تشخیص گفتار، متن به مدل زبانی ارسال می‌شود.
    \item پس از تولید کامل پاسخ متنی، متن به \lr{TTS} ارسال می‌شود.
    \item پس از تولید صوت، پاسخ برای کاربر پخش می‌شود.
\end{enumerate}

این معماری ساده‌تر است و پیاده‌سازی و اشکال‌زدایی آن نیز
راحت‌تر است، اما زمان انتظار کاربر می‌تواند بیشتر باشد.

اگر زمان اجرای هر مرحله به‌ترتیب برابر با

\[
T_{\mathrm{VAD}},
T_{\\mathrm{ASR}},
T_{\\mathrm{LLM}},
T_{\\mathrm{TTS}}
\]

باشد، در یک معماری کاملاً ترتیبی می‌توان تأخیر کلی را به‌طور
تقریبی به شکل زیر بیان کرد:

\begin{equation}
T_{\mathrm{total}}
=
T_{\mathrm{VAD}}
+
T_{\\mathrm{ASR}}
+
T_{\\mathrm{LLM}}
+
T_{\\mathrm{TTS}}
+
T_{\mathrm{overhead}}
\end{equation}

که در آن $T_{\mathrm{overhead}}$ شامل زمان لازم برای انتقال داده،
تبدیل فرمت‌ها و سایر پردازش‌های جانبی است.


% ---------------------------------------------------------
\subsection{پردازش جریانی}

روش دیگر،
\textbf{پردازش جریانی}
\LTRfootnote{Streaming Processing}
است.

در این رویکرد، سامانه برای آغاز پردازش منتظر پایان کامل تمام
مراحل قبلی نمی‌ماند.
برای مثال، \lr{ASR} می‌تواند صوت را به‌صورت قطعات کوتاه دریافت کند
و متن اولیه را به‌تدریج تولید نماید.

همچنین، در برخی معماری‌ها می‌توان تولید پاسخ مدل زبانی را
به‌صورت توکن‌به‌توکن دریافت کرد و بخشی از پاسخ را پیش از
تولید کامل جمله در اختیار \lr{TTS} قرار داد.

به‌صورت مفهومی:

\begin{center}
\lr{Audio} Stream
$\longrightarrow$
Partial \lr{ASR}
$\longrightarrow$
Partial \lr{Text}
$\longrightarrow$
\lr{LLM} Stream
$\longrightarrow$
\lr{TTS} Stream
\end{center}

استفاده از پردازش جریانی می‌تواند زمان شروع پاسخ را کاهش دهد،
اما طراحی سیستم را پیچیده‌تر می‌کند.
در این حالت باید موضوعاتی مانند مدیریت بسته‌های ناقص،
اصلاح خروجی موقت \lr{ASR}، تقسیم پاسخ متنی و هماهنگی میان
تولید متن و تولید صوت در نظر گرفته شوند.

در نسخه اولیه این پروژه، معماری اصلی به‌صورت ترتیبی در نظر
گرفته شده است و امکان استفاده از پردازش جریانی می‌تواند به‌عنوان
یکی از مسیرهای توسعه بعدی بررسی شود.


% ---------------------------------------------------------
\subsection{انتشار خطا در خط لوله}

یکی از چالش‌های مهم معماری زنجیره‌ای،
\textbf{انتشار خطا}
\LTRfootnote{Error Propagation}
است.

فرض شود متن صحیح کاربر برابر باشد با:

\begin{center}
«هوا امروز در تهران چطور است؟»
\end{center}

اگر \lr{ASR} بخشی از جمله را اشتباه تشخیص دهد و خروجی زیر را تولید کند:

\begin{center}
«هوا امروز در تبریز چطور است؟»
\end{center}

مدل زبانی بر اساس همین متن اشتباه پاسخ خواهد داد.
در این حالت، مشکل اصلی در مدل زبانی نیست، بلکه خطای مرحله \lr{ASR}
به مرحله بعد منتقل شده است.

بنابراین خطای کلی سامانه را نمی‌توان تنها با بررسی مستقل
هر مؤلفه توضیح داد.

به‌صورت مفهومی:

\begin{equation}
E_{\mathrm{total}}
=
f(
E_{\\mathrm{ASR}},
E_{\\mathrm{LLM}},
E_{\\mathrm{TTS}},
E_{\mathrm{context}},
E_{\mathrm{preprocess}}
)
\end{equation}

که در آن $E$ نشان‌دهنده خطا یا کاهش کیفیت در هر مرحله است.

این موضوع اهمیت ارزیابی انتهابه‌انتهای سامانه را نشان می‌دهد.


% ---------------------------------------------------------
\subsection{پس‌پردازش متن}

خروجی مدل \lr{ASR} ممکن است نیازمند پردازش تکمیلی باشد.
این پردازش می‌تواند شامل موارد زیر باشد:

\begin{itemize}

    \item نرمال‌سازی فاصله و نیم‌فاصله؛

    \item اصلاح برخی نویسه‌های فارسی و عربی؛

    \item افزودن یا اصلاح علائم نگارشی؛

    \item حذف تکرارهای ناخواسته؛

    \item مدیریت اعداد و علامت‌ها؛

    \item و آماده‌سازی متن برای مدل زبانی.

\end{itemize}

این مرحله اهمیت دارد زیرا مدل زبانی ورودی متنی دریافت می‌کند و
کیفیت این ورودی می‌تواند بر کیفیت پاسخ اثر بگذارد.

در عین حال، پس‌پردازش شدید نیز می‌تواند باعث تغییر نادرست متن
شود.
بنابراین قواعد پس‌پردازش باید با دقت طراحی شوند.


% ---------------------------------------------------------
\subsection{مدیریت بافت مکالمه}

در یک سامانه چندمرحله‌ای، تنها پیام فعلی کاربر برای تولید پاسخ
کافی نیست.
سامانه باید بتواند اطلاعات مهم نوبت‌های قبلی را نیز در اختیار
مدل زبانی قرار دهد.

برای مثال:

\begin{center}
کاربر: یک لپ‌تاپ مناسب برنامه‌نویسی معرفی کن. \\
سامانه: ... \\
کاربر: کدام‌یک سبک‌تر است؟
\end{center}

برای پاسخ به پرسش دوم، مدل باید بداند منظور از «کدام‌یک» به
گزینه‌های معرفی‌شده در پیام قبلی اشاره دارد.

در نتیجه، معماری مکالمه می‌تواند شامل یک مؤلفه
\textbf{مدیریت تاریخچه}
\LTRfootnote{Conversation History Management}
باشد.

به‌صورت ساده:

\begin{equation}
Input_t
=
SystemPrompt
+
History_{t-1}
+
UserMessage_t
\end{equation}

که در آن $History_{t-1}$ شامل بخشی از پیام‌های قبلی مکالمه است.

در پیاده‌سازی عملی باید محدودیت طول بافت مدل زبانی نیز در نظر
گرفته شود.
در صورت طولانی‌شدن مکالمه، ممکن است لازم باشد بخشی از تاریخچه
حذف یا خلاصه شود.


% ---------------------------------------------------------
\subsection{معماری ماژولار}

یکی از اصول مهم طراحی پروژه حاضر،
\textbf{ماژولار بودن}
\LTRfootnote{Modularity}
است.

در یک طراحی ماژولار، هر مؤلفه رابط مشخصی دارد و جزئیات داخلی
آن تا حد امکان از سایر اجزا مستقل است.

برای مثال، ماژول \lr{ASR} را می‌توان به‌صورت مفهومی چنین در نظر گرفت:

\begin{center}
\lr{Audio} Input
$\longrightarrow$
\lr{ASR} Module
$\longrightarrow$
\lr{Text} Output
\end{center}

اگر تمام مدل‌های \lr{ASR} از یک رابط مشترک پیروی کنند، می‌توان
مدل داخلی را تغییر داد بدون آنکه بخش \lr{LLM} یا \lr{TTS} نیازمند
بازنویسی گسترده باشد.

به همین شکل:

\begin{center}
\lr{Text}
$\longrightarrow$
\lr{LLM} Module
$\longrightarrow$
Response \lr{Text}
\end{center}

و:

\begin{center}
\lr{Text}
$\longrightarrow$
\lr{TTS} Module
$\longrightarrow$
\lr{Audio}
\end{center}

این معماری چند مزیت مهم دارد:

\begin{itemize}

    \item امکان جایگزینی آسان مدل‌ها؛

    \item امکان ارزیابی مستقل هر مؤلفه؛

    \item کاهش وابستگی میان بخش‌های نرم‌افزار؛

    \item ساده‌ترشدن اشکال‌زدایی؛

    \item و امکان توسعه آینده سامانه.

\end{itemize}

این ویژگی برای پروژه حاضر اهمیت زیادی دارد، زیرا چند مدل
مختلف برای هر یک از اجزای \lr{ASR}، \lr{LLM} و \lr{TTS} مورد آزمایش
قرار گرفته‌اند.


% ---------------------------------------------------------
\subsection{مدیریت خطا}

در یک سامانه یکپارچه، ممکن است هر یک از مؤلفه‌ها با خطا مواجه
شوند.

برای مثال:

\begin{itemize}

    \item میکروفون ممکن است صوت مناسبی دریافت نکند؛

    \item VAD ممکن است گفتاری تشخیص ندهد؛

    \item \lr{ASR} ممکن است متن خالی یا نامعتبر تولید کند؛

    \item مدل زبانی ممکن است با کمبود حافظه مواجه شود؛

    \item \lr{TTS} ممکن است نتواند متن خاصی را پردازش کند؛

    \item یا فایل صوتی خروجی به‌درستی تولید نشود.

\end{itemize}

اگر این خطاها مدیریت نشوند، یک مشکل کوچک می‌تواند باعث توقف
کامل سامانه شود.

بنابراین معماری باید دارای سازوکارهایی برای تشخیص خطا،
ثبت گزارش و بازیابی مناسب باشد.

در سطح مفهومی:

\begin{center}
Module Execution
$\longrightarrow$
Success
$\longrightarrow$
Next Module
\end{center}

یا:

\begin{center}
Module Execution
$\longrightarrow$
Failure
$\longrightarrow$
Error Handler
\end{center}

استفاده از ثبت رخدادها
\LTRfootnote{Logging}
نیز می‌تواند در تشخیص منبع خطا و اندازه‌گیری زمان اجرای
مؤلفه‌ها مفید باشد.


% ---------------------------------------------------------
\subsection{مصرف هم‌زمان منابع}

یکی از تفاوت‌های مهم ارزیابی مستقل مدل‌ها با اجرای کامل سامانه
این است که در حالت یکپارچه، چند مدل ممکن است روی یک سخت‌افزار
مشترک اجرا شوند.

برای مثال، اگر:

\[
M_{\\mathrm{ASR}},
M_{\\mathrm{LLM}},
M_{\\mathrm{TTS}}
\]

به‌ترتیب حافظه مورد نیاز هر مدل باشند، مقدار حافظه لازم در
یک اجرای ساده می‌تواند تقریباً برابر باشد با:

\begin{equation}
M_{\mathrm{total}}
\approx
M_{\\mathrm{ASR}}
+
M_{\\mathrm{LLM}}
+
M_{\\mathrm{TTS}}
+
M_{\mathrm{runtime}}
\end{equation}

در عمل مقدار دقیق به نحوه بارگذاری مدل‌ها، آزادسازی حافظه،
دقت عددی و چارچوب مورد استفاده وابسته است.

بنابراین ممکن است مدلی در ارزیابی مستقل قابل اجرا باشد، اما
هنگام بارگذاری هم‌زمان با سایر مدل‌ها محدودیت حافظه ایجاد کند.

این مسئله یکی از دلایل اصلی اهمیت انتخاب مدل بر اساس توازن
میان کیفیت و مصرف منابع است.


% ---------------------------------------------------------
\subsection{تأخیر انتهابه‌انتها}

معیار مهم دیگر در یک چت‌بات صوتی،
\textbf{تأخیر انتهابه‌انتها}
\LTRfootnote{End-to-End \lr{Latency}}
است.

این معیار نشان می‌دهد از زمانی که کاربر صحبت خود را پایان
می‌دهد تا زمانی که پاسخ سامانه آماده یا پخش می‌شود چه مدت
زمان سپری شده است.

یک مدل ساده برای این تأخیر عبارت است از:

\begin{equation}
T_{\mathrm{E2E}}
=
T_{\mathrm{endpoint}}
+
T_{\\mathrm{ASR}}
+
T_{\mathrm{text}}
+
T_{\\mathrm{LLM}}
+
T_{\\mathrm{TTS}}
+
T_{\mathrm{I/O}}
\end{equation}

که در آن:

\begin{itemize}

    \item $T_{\mathrm{endpoint}}$ زمان تشخیص پایان گفتار؛
    \item $T_{\\mathrm{ASR}}$ زمان تبدیل صوت به متن؛
    \item $T_{\mathrm{text}}$ زمان پردازش متنی؛
    \item $T_{\\mathrm{LLM}}$ زمان تولید پاسخ؛
    \item $T_{\\mathrm{TTS}}$ زمان تولید صوت؛
    \item و $T_{\mathrm{I/O}}$ زمان مرتبط با ورودی و خروجی است.

\end{itemize}

این تجزیه کمک می‌کند مشخص شود کدام مؤلفه بیشترین سهم را در
تأخیر کلی دارد.

برای مثال، اگر مدل \lr{ASR} تنها 0.1 ثانیه زمان نیاز داشته باشد
اما مدل زبانی 3 ثانیه برای تولید پاسخ مصرف کند، بهینه‌سازی
بیشتر \lr{ASR} تأثیر محدودی بر تجربه کلی کاربر خواهد داشت.

ازاین‌رو، در فصل آزمایش‌ها علاوه بر زمان مستقل هر مدل،
اندازه‌گیری تأخیر انتهابه‌انتهای خط لوله نیز اهمیت خواهد داشت.


% ---------------------------------------------------------
\subsection{معماری پایه پروژه حاضر}

بر اساس مفاهیم مطرح‌شده، معماری پایه پروژه حاضر را می‌توان
به شکل زیر خلاصه کرد:

\begin{center}
User \lr{Speech}
$\longrightarrow$
\lr{Audio} Input
$\longrightarrow$
\lr{ASR}
$\longrightarrow$
\lr{Text}
$\longrightarrow$
\lr{LLM}
$\longrightarrow$
Response \lr{Text}
$\longrightarrow$
\lr{TTS}
$\longrightarrow$
Response \lr{Audio}
\end{center}

نسخه پایه این معماری به‌صورت محلی پیاده‌سازی شده و ارتباط
اصلی میان سه ماژول
\lr{ASR}، \lr{LLM} و \lr{TTS}
برقرار شده است.

در ادامه توسعه پروژه، مؤلفه‌هایی مانند VAD، پردازش نویز،
مدیریت دقیق‌تر بافت مکالمه، پس‌پردازش متن، مدیریت خطا و
اندازه‌گیری دقیق تأخیر به این معماری اضافه خواهند شد.

جزئیات طراحی این معماری و نحوه ارتباط دقیق ماژول‌ها در
فصل توصیف پروژه ارائه می‌شود، در حالی که جزئیات پیاده‌سازی
نرم‌افزاری در فصل پیاده‌سازی تشریح خواهد شد.
% =========================================================
% 2-9
% =========================================================
\section{جمع‌بندی}

در این فصل، مفاهیم نظری مورد نیاز برای طراحی و تحلیل یک
چت‌بات صوتی فارسی معرفی شدند.

ابتدا مبانی پردازش گفتار شامل نمونه‌برداری، قاب‌بندی،
نمایش‌های حوزه زمان و فرکانس، طیف‌نگاشت، مقیاس Mel و تأثیر
نویز بر سیگنال گفتار بررسی شد.
سپس اصول تشخیص خودکار گفتار و معیارهای اصلی ارزیابی آن شامل
\lr{WER}، \lr{CER}، \lr{RTF} و زمان تأخیر تشریح شدند.

در ادامه، معماری‌های تشخیص گفتار مورد استفاده در پروژه شامل
\lr{wav2vec 2.0/XLSR}، \lr{Whisper}، \lr{Vosk/Kaldi}
و \lr{FastConformer Hybrid} معرفی شدند و تفاوت مفهومی
رویکردهای \lr{CTC}، رمزگشایی خودبازگشتی و
\lr{RNN-T} مورد بررسی قرار گرفت.

پس از آن، مفاهیم پایه مدل‌های زبانی بزرگ شامل معماری
\lr{Transformer}، توجه، مدل‌های خودبازگشتی، توکن‌سازی،
تنظیم دستورمحور، مدیریت بافت مکالمه و اجرای محلی مدل‌ها
بررسی شدند.
همچنین معیارهایی مانند زمان پاسخ، نرخ تولید توکن و مصرف
حافظه برای ارزیابی عملی مدل‌های زبانی معرفی شدند.

در بخش تبدیل متن به گفتار، ساختار کلی سامانه‌های \lr{TTS}،
معماری \lr{VITS}، مدل‌های \lr{MMS} و سامانه \lr{Piper}
بررسی شدند.
معیارهای ارزیابی کیفیت و کارایی شامل \lr{MOS}، \lr{NISQA}،
\lr{WER}، \lr{CER} و \lr{RTF} نیز تشریح شدند.

در ادامه، نقش تشخیص فعالیت صوتی و کاهش نویز در آماده‌سازی
ورودی صوتی بررسی شد.
نشان داده شد که این مؤلفه‌ها علاوه بر دقت تشخیص گفتار،
می‌توانند بر تأخیر کلی سامانه نیز اثرگذار باشند.

در نهایت، معماری عمومی سامانه‌های گفت‌وگوی صوتی،
پردازش ترتیبی و جریانی، انتشار خطا، مدیریت بافت،
طراحی ماژولار، مصرف منابع و تأخیر انتهابه‌انتها مورد بررسی
قرار گرفت.

مفاهیم ارائه‌شده در این فصل، مبنای نظری لازم برای بررسی
پژوهش‌های مرتبط در فصل بعد و سپس تشریح معماری و پیاده‌سازی
سامانه پیشنهادی را فراهم می‌کنند.
